% File: Volume-I-Foundations/main-vol1.tex

\documentclass[11pt]{book}

% Relative paths assume this file is in Volume-I-Foundations/
% File: shared/macros.tex
% Global macros and packages shared across all three volumes.

% ------------------------------------------------------------
% Encoding, fonts, layout
% ------------------------------------------------------------
\usepackage[T1]{fontenc}
\usepackage[utf8]{inputenc}
\usepackage{lmodern}

\usepackage[a4paper,margin=1in]{geometry}
\usepackage{setspace}
\onehalfspacing

\usepackage{microtype}

% ------------------------------------------------------------
% Math packages
% ------------------------------------------------------------
\usepackage{amsmath,amssymb,amsfonts,mathtools}
\usepackage{bm}
\usepackage{mathrsfs}

% For diagrams (internally; no ASCII diagrams in text)
\usepackage{tikz-cd}

% ------------------------------------------------------------
% Theorem environments (delegated to theoremstyles.tex)
% ------------------------------------------------------------
% File: shared/theoremstyles.tex
% Centralized theorem styles and numbering scheme.

\usepackage{amsthm}

% Chapter-based numbering
\numberwithin{equation}{chapter}

\theoremstyle{plain}
\newtheorem{theorem}{Theorem}[chapter]
\newtheorem{proposition}[theorem]{Proposition}
\newtheorem{lemma}[theorem]{Lemma}
\newtheorem{corollary}[theorem]{Corollary}
\newtheorem{conjecture_env}[theorem]{Conjecture}

\theoremstyle{definition}
\newtheorem{definition}[theorem]{Definition}
\newtheorem{example}[theorem]{Example}

\theoremstyle{remark}
\newtheorem{remark}[theorem]{Remark}

% Environment for explicitly marked open problems
\newtheorem{openproblem}[theorem]{Open Problem}



% ------------------------------------------------------------
% Hyperlinks and references
% ------------------------------------------------------------
\usepackage{hyperref}
\hypersetup{
  colorlinks=true,
  linkcolor=blue,
  citecolor=blue,
  urlcolor=blue
}

\usepackage[nameinlink,capitalize]{cleveref}

% ------------------------------------------------------------
% Bibliography (volume-specific files will configure biblatex)
% ------------------------------------------------------------
% Each main-*.tex will load biblatex with its own .bib file.

% ------------------------------------------------------------
% Common macros
% ------------------------------------------------------------

% Sets and categories
\newcommand{\R}{\mathbb{R}}
\newcommand{\C}{\mathbb{C}}
\newcommand{\Z}{\mathbb{Z}}
\newcommand{\N}{\mathbb{N}}

\newcommand{\Cat}{\mathsf{Cat}}
\newcommand{\sSet}{\mathsf{sSet}}
\newcommand{\Top}{\mathsf{Top}}

% RSVP core fields
\newcommand{\PhiField}{\Phi}
\newcommand{\vField}{\bm{v}}
\newcommand{\SField}{S}

% Differential geometry
\newcommand{\dd}{\mathrm{d}}
\newcommand{\grad}{\nabla}
\newcommand{\Div}{\operatorname{div}}
\newcommand{\Lap}{\Delta}
\newcommand{\Lie}{\mathcal{L}}

% Moduli stacks and derived notions
\newcommand{\Plenum}{\mathsf{Plenum}}
\newcommand{\RPlenum}{\mathbb{R}\mathsf{Plenum}}
\newcommand{\Map}{\operatorname{Map}}
\newcommand{\Spec}{\operatorname{Spec}}
\newcommand{\colim}{\operatorname{colim}}
\newcommand{\hocolim}{\operatorname{hocolim}}

% Symplectic / BV / AKSZ notation
\newcommand{\Tstarn}[1]{T^{\ast}[#1]}
\newcommand{\Tstar}[1]{T^{\ast}#1}
\newcommand{\SympShifted}[1]{#1\text{--shifted symplectic}}
\newcommand{\BV}{\mathrm{BV}}
\newcommand{\AKSZ}{\mathrm{AKSZ}}
\newcommand{\CM}{\mathrm{CM}} % Classical master equation

% SpherePop
\newcommand{\SpherePop}{\mathsf{SpherePop}}
\newcommand{\Collapse}{\kappa}

% Environments for “intuition”, “warning”, etc.
\newcommand{\Intuition}{\paragraph{Intuition.}}
\newcommand{\Warning}{\paragraph{Warning.}}
\newcommand{\HistoricalNote}{\paragraph{Historical note.}}

% ------------------------------------------------------------
% Volume–wide metadata helpers
% ------------------------------------------------------------
\newcommand{\rsvptitle}{Relativistic Scalar--Vector Plenum}
\newcommand{\rsvpabbr}{RSVP}



% ------------------------------------------------------------
% Bibliography setup (volume-specific)
% ------------------------------------------------------------
\usepackage[backend=biber,style=alphabetic,maxbibnames=10]{biblatex}
\addbibresource{bibliography-vol1.bib}

% ------------------------------------------------------------
% Title metadata
% ------------------------------------------------------------
\title{\rsvptitle:\\
Mathematical Foundations of Lamphrodynamic Field Theory}
\author{Flyxion}
\date{\today}

\begin{document}

\frontmatter
\maketitle

\chapter*{Preface}

This volume is the first of a three--volume monograph developing the
Relativistic Scalar--Vector Plenum (\rsvpabbr{}) as a rigorous mathematical
framework. Its primary audience is mathematicians and theoretical physicists
with a background in differential geometry, category theory, and quantum field
theory. A secondary aim is pedagogical: the material is written so that
physicists willing to learn derived geometry can follow, provided they are
comfortable with homological algebra and gauge theory.

The central goal of Volume~I is to construct the lamphrodynamic field theory
underlying \rsvpabbr{}, to formulate it in the language of derived algebraic
geometry and shifted symplectic structures, and to establish fundamental
analytic results (well--posedness, energy estimates, regularity). Later volumes
will build on this foundation to explore cosmology, emergent gravity, geometric
computation, and philosophical implications.

We have tried to adhere to the following design principles:
\begin{itemize}
  \item \emph{Rigor first in foundations.} The core constructions (derived
  stacks, shifted symplectic forms, AKSZ action) are presented with as much
  mathematical rigor as possible.

  \item \emph{Motivation before formalism.} Each chapter begins with an
  informal explanation of why the subsequent technical machinery is needed.

  \item \emph{Explicit status of claims.} Theorems are proved or referenced;
  conjectures and open problems are clearly labeled and accompanied by
  evidence or heuristic arguments.
\end{itemize}

Volume~I is structured in four parts. Part~I reviews the necessary background
in derived geometry and shifted symplectic structures, oriented toward
physicists. Part~II introduces the lamphrodynamic plenum and constructs the
relevant derived moduli stacks and AKSZ action. Part~III develops the analytic
foundations of the lamphrodynamic PDEs. Part~IV investigates geometric and
topological properties, including entropic singularities and symmetries.

Subsequent volumes will not repeat this material, but will refer back to it
frequently.

\vspace{1em}
\noindent\textbf{Acknowledgements.} 
The author thanks numerous interlocutors --- human and artificial --- for
many conversations that shaped this framework. Any errors are, of course,
the author’s responsibility.

\tableofcontents

\chapter*{Notation and Conventions}

Unless otherwise stated, we adopt the following conventions:
\begin{itemize}
  \item $(M,g)$ denotes a smooth oriented Lorentzian or Riemannian manifold,
  with metric $g$ of signature $(-,+,\dots,+)$ in the Lorentzian case.

  \item Greek indices $\mu,\nu,\dots$ range over spacetime coordinates;
  Latin indices $i,j,\dots$ are reserved for spatial coordinates when a
  splitting is chosen. We freely use the abstract index notation of Penrose.

  \item The Levi--Civita connection of $g$ is denoted by $\nabla$ and the
  associated Laplace--Beltrami operator by $\Delta = \nabla_a\nabla^a$.

  \item The symbol $\Map(X,Y)$ denotes a (derived) mapping stack when $X$ and
  $Y$ are derived stacks, and an ordinary mapping space when they are topological
  or smooth manifolds; the intended sense is clear from context.

  \item All tensor contractions use the Einstein summation convention.

  \item The core lamphrodynamic fields of \rsvpabbr{} are denoted
  $(\PhiField,\vField,\SField)$: a scalar potential, a vector flow, and an
  entropy density. We often group them into a single configuration field
  $\Psi := (\PhiField,\vField,\SField)$.
\end{itemize}

We recommend that readers unfamiliar with derived stacks or shifted symplectic
geometry skim \cref{chap:intro,chap:inf-cat} first, and then consult
\cref{chap:shifted-symplectic,chap:aksz} as needed when these notions appear
later.

\mainmatter

% ------------------------------------------------------------
% Part I: Derived Algebraic Geometry for Physicists
% ------------------------------------------------------------
\part{Derived Algebraic Geometry for Physicists}

% File: Volume-I-Foundations/chapters/ch01-intro.tex

\chapter{Introduction and Motivation}
\label{chap:intro}

The Relativistic Scalar--Vector Plenum (\rsvpabbr{}) is a proposal for a
unified field--theoretic framework in which gravity, cosmology, and certain
aspects of information and agency arise from the dynamics of a triplet of
fields
\[
  (\PhiField,\vField,\SField)
\]
living on a spacetime manifold $(M,g)$ or, more generally, on a derived
geometric object that refines $(M,g)$. The scalar field $\PhiField$ encodes a
form of lamphrodynamic potential, $\vField$ represents a flux or flow of
negentropic influence, and $\SField$ is an entropy density that couples to
geometry and matter.

At a heuristic level, \rsvpabbr{} treats spacetime not as a static stage on
which fields live, but as a \emph{plenum} whose local structure is determined
by the coupled evolution of $(\PhiField,\vField,\SField)$. Geometrical
quantities such as curvature, torsion, and metric relations are then derived
from the lamphrodynamic configuration rather than posited \emph{a priori}.

This introductory chapter explains why such a framework naturally leads us to
derived algebraic geometry and shifted symplectic structures, and sets the
stage for the rigorous constructions that follow.

%-----------------------------------------------------------------
\section{From Classical Fields to Derived Geometry}
%-----------------------------------------------------------------

In classical field theory, a field configuration is a section of some
bundle over spacetime $M$, and the space of fields is modeled as an infinite
dimensional manifold or Fréchet space. Functional integrals, symmetries, and
constraints are handled with varying degrees of rigor depending on context.

However, many modern developments --- such as moduli spaces of solutions to
gauge equations, intersection theory on spaces of instantons, and the BV
formalism for gauge theories --- reveal that naive smooth manifolds are
insufficient. One encounters:
\begin{itemize}
  \item Singularities when solutions collide or degenerate.
  \item Higher automorphisms from gauge symmetries.
  \item Obstructions and derived intersections.
\end{itemize}

Derived algebraic geometry was developed precisely to give these moduli spaces
a well--behaved home. Instead of regarding a moduli space as a naive quotient
or set of solutions, one encodes it as a derived stack equipped with additional
structure (cotangent complex, obstruction theory, shifted symplectic form).

\Intuition
Even if the base spacetime $M$ is an ordinary smooth manifold, the \emph{space
of fields} is almost never a manifold in the classical sense. Gauge symmetry,
constraints, and singularities force us to work in a homotopical setting.
Derived stacks and $\infty$--categories are the natural language for this.

In \rsvpabbr{}, lamphrodynamic fields are subject to:
\begin{itemize}
  \item Gauge transformations mixing $\PhiField$ and $\vField$,
  \item Constraints relating $\SField$ to curvature and matter content,
  \item Entropic singularities where $\SField$ becomes non--smooth or diverges.
\end{itemize}
These features suggest that the moduli of lamphrodynamic configurations should
be modeled as a derived stack, which we denote $\Plenum$. A significant portion
of Volume~I is devoted to constructing $\Plenum$, equipping it with a
$(-1)$--shifted symplectic structure, and understanding its tangent and
obstruction theory.

%-----------------------------------------------------------------
\section{The Lamphrodynamic Triplet}
%-----------------------------------------------------------------

We give a preliminary, heuristic description of the fields before formalizing
them in \cref{chap:lamphrodynamic-axioms,chap:plenum-stack}.

\begin{itemize}
  \item $\PhiField$ is a real--valued scalar field on $M$ (or a section of
  some line bundle) representing a lamphrodynamic potential. Regions of high
  $\PhiField$ can be thought of as reservoirs of ``available negentropy''.

  \item $\vField$ is a vector field on $M$, or more generally a section of
  $TM\otimes E$ for some auxiliary bundle $E$, encoding the directed transport
  of lamphrodynamic influence. It plays a role analogous to a current or
  velocity field.

  \item $\SField$ is a scalar entropy density that couples to both geometry and
  matter fields. In contrast to thermodynamic entropy defined on macrostates,
  $\SField$ is intended as a local field that participates dynamically in the
  evolution equations.
\end{itemize}

The basic lamphrodynamic field equations are schematically of the form
\begin{align}
  \partial_t \PhiField &=
    \mathcal{D}_\Phi \Delta \PhiField
    + \mathcal{F}_\Phi(\PhiField,\vField,\SField,g), \label{eq:intro-phi}\\
  \partial_t \SField &=
    \mathcal{D}_S \Delta \SField
    + \mathcal{F}_S(\PhiField,\vField,\SField,g), \label{eq:intro-S}\\
  \partial_t \vField &=
    \mathcal{D}_v \Delta \vField
    + \mathcal{F}_v(\PhiField,\vField,\SField,g), \label{eq:intro-v}
\end{align}
where $\mathcal{D}_\Phi,\mathcal{D}_S,\mathcal{D}_v$ are diffusion constants
and the nonlinear terms $\mathcal{F}_{\Phi,S,v}$ encode lamphrodynamic
interactions and coupling to curvature. Precise versions of these equations
will be derived from an action principle in \cref{chap:aksz-rsvp,chap:variational-hamiltonian}.

\Warning
Equations~\eqref{eq:intro-phi}--\eqref{eq:intro-v} are merely schematic at this
stage. They suggest that \rsvpabbr{} is, at its analytic core, a system of
parabolic or mixed--type PDEs; however, the actual lamphrodynamic equations
will include nonlocal and higher--order terms arising from the BV--AKSZ
construction.

%-----------------------------------------------------------------
\section{Why a $(-1)$--Shifted Symplectic Structure?}
%-----------------------------------------------------------------

In classical Hamiltonian mechanics, the phase space of a system is a symplectic
manifold $(X,\omega)$ and the dynamics are encoded by a Hamiltonian function
$H$ via the equation
\[
  \iota_{X_H}\omega = \dd H.
\]
For gauge theories and field theories, the BV formalism generalizes this
picture to a graded setting: one replaces $(X,\omega)$ by a graded manifold
equipped with an odd symplectic form, and the Hamiltonian $H$ by a master
action $S$ satisfying $\{S,S\}=0$.

In derived geometry, these graded symplectic manifolds are captured by
\emph{shifted symplectic structures} on derived stacks. A $k$--shifted
symplectic structure is, roughly, a closed, nondegenerate $2$--form of
cohomological degree $k$ on the derived stack. The case $k=-1$ is particularly
important for the BV formalism: it corresponds to odd symplectic forms and is
tailored to encode gauge theories.

\begin{definition}
Let $\mathcal{X}$ be a derived Artin stack over a field of characteristic
zero. A \emph{$(-1)$--shifted symplectic structure} on $\mathcal{X}$ is a
closed $2$--form
\[
  \omega \in \mathcal{A}^{2,cl}(\mathcal{X})[-1]
\]
whose underlying pairing on the cotangent complex is nondegenerate in the
derived sense.
\end{definition}

The key insight for \rsvpabbr{} is that the moduli stack $\Plenum$ of
lamphrodynamic configurations should carry such a $(-1)$--shifted symplectic
structure, with the master action functional $S_{\rsvpabbr}$ playing the role
of a Hamiltonian satisfying the classical master equation
\[
  \{S_{\rsvpabbr},S_{\rsvpabbr}\} = 0.
\]
Once this structure is constructed, the AKSZ formalism provides a natural way
to build $S_{\rsvpabbr}$ from geometric data associated to $\Plenum$ and the
spacetime source manifold.

\Intuition
The choice of $k=-1$ is not arbitrary: lamphrodynamic fields are subject to
gauge redundancies and constraints, so their moduli naturally form a BV phase
space rather than a naive symplectic manifold. This forces the symplectic
structure to be of odd degree, hence $k=-1$.

A central technical goal of Volume~I is to:
\begin{itemize}
  \item Construct $\Plenum$ as a derived Artin stack.
  \item Equip $\Plenum$ with a $(-1)$--shifted symplectic structure $\omega$.
  \item Show that $\omega$ can be expressed in terms of $(\PhiField,\vField,\SField)$
  and curvature data, and that it is closed and nondegenerate.
\end{itemize}
These results will be established in \cref{chap:plenum-stack,chap:aksz-rsvp}.

%-----------------------------------------------------------------
\section{Analytic Questions: Well–Posedness and Regularity}
%-----------------------------------------------------------------

The derived geometric formulation of \rsvpabbr{} provides a powerful conceptual
framework, but physics also demands that the associated PDEs be analytically
well--behaved. Concretely, we must address:

\begin{enumerate}
  \item \emph{Well--posedness.} Given initial data $(\PhiField_0,\vField_0,\SField_0)$,
  does there exist a unique solution $(\PhiField,\vField,\SField)$, at least
  for short time, in appropriate Sobolev spaces?

  \item \emph{Energy estimates.} Are there a priori bounds controlling the
  growth of norms of the fields, and do these bounds depend continuously on
  the initial data?

  \item \emph{Regularity and singularities.} Under what conditions can
  singularities form in finite time, and how are such singularities related to
  entropic phenomena in \rsvpabbr{}?

  \item \emph{Long--time behavior.} Do solutions converge to lamphrodynamic
  equilibria or attractors, and can one classify these asymptotic states?
\end{enumerate}

These questions are addressed in Part~III of this volume. We will prove
global well--posedness for linearized lamphrodynamic equations and establish
energy estimates and local well--posedness for certain nonlinear regimes.
Global existence for the fully nonlinear system remains conjectural; we will
state precise open problems and present partial evidence.

%-----------------------------------------------------------------
\section{Organization of the Three Volumes}
%-----------------------------------------------------------------

For orientation, we briefly summarize the roles of the three volumes.

\subsection*{Volume I: Mathematical Foundations}

This volume develops:
\begin{itemize}
  \item The background in derived algebraic geometry and shifted symplectic
  structures needed to formulate \rsvpabbr{} rigorously.
  \item The construction of the lamphrodynamic moduli stack $\Plenum$ and its
  $(-1)$--shifted symplectic structure.
  \item The AKSZ--BV action for lamphrodynamic fields and its classical master
  equation.
  \item Fundamental analytic results on well--posedness, energy estimates, and
  regularity of the lamphrodynamic PDEs.
\end{itemize}

\subsection*{Volume II: Physical Applications}

The second volume applies this foundation to physics:
\begin{itemize}
  \item Derivation of modified Einstein equations with an explicit lamphrodynamic
  stress--energy contribution $\Theta_{ab}$.
  \item Development of a cosmological model in which redshift, structure
  formation, and cosmic microwave background anisotropies are governed by
  lamphrodynamic fields rather than pure metric expansion.
  \item Detailed comparison with observational data (supernovae, BAO, Planck
  CMB measurements) and identification of potential tests that distinguish
  \rsvpabbr{} from $\Lambda$CDM.
\end{itemize}

\subsection*{Volume III: Extensions and Interpretations}

The third volume explores:
\begin{itemize}
  \item The SpherePop calculus as a geometric model of computation, including
  categorical and homotopical formalization and a proof of Turing completeness.
  \item L--systems, semantic fields, and distributed computation (CRDTs) as
  instances of lamphrodynamic or related field theories.
  \item Entropy--weighted spectral geometry and ``cymatic'' resonance modes.
  \item Philosophical and creative dimensions, including the interpretation of
  \rsvpabbr{} within structural realism and process philosophy.
\end{itemize}

%-----------------------------------------------------------------
\section{What is Proven, What is Conjectural}
%-----------------------------------------------------------------

Given the ambitious scope of \rsvpabbr{}, it is important to distinguish
carefully between results that are established in this monograph and those
that remain conjectural.

\begin{itemize}
  \item \emph{Proven in Volume I:} We provide complete proofs for the
  construction of $\Plenum$ as a derived Artin stack, for the existence of a
  $(-1)$--shifted symplectic form on $\Plenum$ under specified conditions, and
  for well--posedness results in the linear and certain nonlinear lamphrodynamic
  regimes.

  \item \emph{Proven in Volume II:} We derive modified Einstein equations from
  the lamphrodynamic action, obtain concrete formulas for cosmological
  observables (including a redshift function $f(\SField,\vField)$), and
  perform quantitative comparisons with existing data. The degree of
  observational agreement is investigated but not treated as a fully
  conclusive fit.

  \item \emph{Proven in Volume III:} We rigorously formalize SpherePop as a
  homotopy--theoretic model of computation, prove that merge operations are
  homotopy colimits, and establish Turing completeness via explicit encodings.

  \item \emph{Conjectural:} Global existence of solutions for the full
  nonlinear lamphrodynamic system; completeness of the cosmological model as a
  full replacement for $\Lambda$CDM; certain speculative connections between
  lamphrodynamic spectral modes and cognition or consciousness.
\end{itemize}

These conjectural aspects are clearly labeled as such, and in each case we
provide reasons for their plausibility and suggest directions for future
research.

%-----------------------------------------------------------------
\section{How to Read this Volume}
%-----------------------------------------------------------------

Different readers may wish to approach the material in different orders.

\begin{itemize}
  \item \textbf{Mathematical physicists} familiar with homological algebra and
  derived geometry may skim Part~I and focus on Parts~II and~III, where the
  lamphrodynamic stack and PDE analysis are developed.

  \item \textbf{Physicists new to derived geometry} are encouraged to read
  \cref{chap:inf-cat,chap:shifted-symplectic} carefully, referring to
  appendices as needed, before tackling the AKSZ construction in
  \cref{chap:aksz-bv,chap:aksz-rsvp}.

  \item \textbf{Analysts} interested in PDE aspects may wish to jump directly
  to Part~III, using the early chapters primarily for motivation and notation.
\end{itemize}

Throughout, we have tried to separate background material, core constructions,
and speculative interpretations. The reader is invited to treat this volume as
a foundation that can be built upon independently of whether they ultimately
accept the full physical picture proposed by \rsvpabbr{}.

\bigskip
In the next chapter, we begin our technical journey by reviewing the basic
language of model categories, $\infty$--categories, and derived stacks, with
an eye toward their eventual application to lamphrodynamic field theory.


% File: Volume-I-Foundations/chapters/ch02-infinity-categories.tex

\chapter{\texorpdfstring{$\infty$}{Infinity}--Categories and Homotopy Theory}
\label{chap:inf-cat}

This chapter reviews the homotopical and $\infty$--categorical structures
needed to formulate derived geometry and, ultimately, the lamphrodynamic moduli
stack $\Plenum$ developed in later chapters. We assume that the reader is
familiar with basic category theory and classical homotopy theory. Our goal is
not to give a comprehensive treatment, but to present a minimal collection of
definitions, constructions, and results that will be invoked in the sequel.

Throughout, we work over a field of characteristic zero (typically $\R$ or
$\C$) unless otherwise specified. References include \cite{Lurie2009HTT},
\cite{Lurie2017HA}, and \cite{ToenVezzosi2008}.

%---------------------------------------------------------------
\section{Model Categories and Homotopical Structures}
%---------------------------------------------------------------

The motivation for $\infty$--categories begins with the desire to regard
homotopy equivalent objects as ``the same'' while retaining control over
higher homotopies. Classical category theory does not remember higher
homotopies; instead, it collapses homotopies to identities. In homotopy
theory, however, we need to keep track of:
\begin{itemize}
  \item homotopies between maps,
  \item homotopies between homotopies,
  \item and so on, ad infinitum.
\end{itemize}

A traditional approach is provided by \emph{model categories}, which endow a
category with distinguished classes of weak equivalences, fibrations, and
cofibrations satisfying Quillen's axioms \cite{Quillen1967}. The basic example
is the model structure on topological spaces, where weak equivalences are
weak homotopy equivalences, fibrations are Serre fibrations, and cofibrations
are retracts of relative CW--complex inclusions.

While model categories provide a powerful toolkit, they are ultimately an
\emph{encoding} of $\infty$--categorical structures, and the language of
$\infty$--categories gives a more conceptual, coordinate--free approach.

%---------------------------------------------------------------
\section{\texorpdfstring{$\infty$}{Infinity}--Categories}
%---------------------------------------------------------------

An $\infty$--category is a category enriched in spaces rather than sets, or,
equivalently, a homotopical object that remembers all higher morphisms. A
concrete model is given by \emph{quasi--categories} (simplicial sets
satisfying inner horn conditions) introduced by Boardman and Vogt and
developed extensively by Joyal and Lurie.

\begin{definition}
A \emph{quasi--category} is a simplicial set $C$ such that every inner
horn $\Lambda^n_i \to C$ admits a filler $\Delta^n\to C$ for $0 < i < n$.
\end{definition}

This condition generalizes the Kan condition for simplicial sets
corresponding to $\infty$--groupoids (spaces). Every topological space has a
singular complex which is a Kan complex; every Kan complex is a model for an
$\infty$--groupoid. The quasi--category condition weakens Kan filling so that
objects need not have all inverses, thus modeling $\infty$--categories rather
than $\infty$--groupoids.

\Intuition
Ordinary categories forget higher homotopies between morphisms. Quasi--categories
remember not only morphisms but also homotopies and homotopies between homotopies,
all the way up. This makes them ideal for encoding moduli spaces, where higher
automorphisms naturally appear.

%---------------------------------------------------------------
\section{$\infty$--Topoi and Sheaves of Spaces}
%---------------------------------------------------------------

Derived geometry requires a setting where ``sheaves of homotopical data'' can
be treated systematically. This motivates the notion of an $\infty$--topos:
a homotopical enhancement of Grothendieck topoi.

\begin{definition}
An \emph{$\infty$--topos} is an $\infty$--category that is equivalent to an
$\infty$--category of sheaves of spaces on a Grothendieck site. Equivalently,
it is a presentable $\infty$--category satisfying descent for hypercovers
\cite{Lurie2009HTT}.
\end{definition}

Classical sheaf theory associates sets or vector spaces to open subsets of a
manifold. An $\infty$--topos associates \emph{spaces} (i.e.\ $\infty$--groupoids)
to open subsets, thereby encoding higher homotopical information. Derived
stacks are then defined as (higher) sheaves of $\infty$--groupoids satisfying
appropriate representability and local geometricity conditions.

\begin{remark}
The $\infty$--topos of sheaves of spaces on the site $(\mathrm{Sch},\mathrm{Zariski})$
is denoted ${\rm Shv}_{\infty}(\mathrm{Sch})$. Derived algebraic geometry works
mostly inside suitable subcategories of this $\infty$--topos.
\end{remark}

%---------------------------------------------------------------
\section{Derived Stacks}
%---------------------------------------------------------------

Derived stacks were developed by Toën–Vezzosi, Lurie, and others to handle
moduli problems with obstructions, singularities, and nontrivial automorphism
groups. Roughly speaking, derived stacks are functors from derived rings to
$\infty$--groupoids satisfying descent and representability conditions.

\begin{definition}
A \emph{derived stack} is a functor
\[
  F : \mathrm{dAlg}^{op} \longrightarrow \infty\text{--}\mathrm{Gpd}
\]
satisfying $\infty$--geometricity and descent conditions.
\end{definition}

Derived stacks generalize classical schemes and algebraic stacks, adding:
\begin{itemize}
  \item derived intersections,
  \item cotangent complexes,
  \item obstruction theories,
  \item and higher automorphisms.
\end{itemize}

\Intuition
In physics, moduli spaces of solutions to gauge equations (e.g.\ Yang--Mills,
GR) often have singularities and higher gauge symmetries. Derived stacks treat
these automatically, removing the need for ad hoc patching arguments.

In \rsvpabbr{}, the moduli stack $\Plenum$ of lamphrodynamic configurations
will be constructed as a derived Artin stack equipped with a $(-1)$--shifted
symplectic structure. This stack encodes field configurations, gauge orbits,
and entropic singularities in a unified homotopical object.

%---------------------------------------------------------------
\section{Tangent Complexes and Obstruction Theory}
%---------------------------------------------------------------

Classical manifolds possess tangent spaces at each point, controlling
infinitesimal deformations. Derived stacks generalize this to a \emph{tangent
complex}, which is typically a chain complex, possibly with cohomology in
multiple degrees.

\begin{definition}
Let $x\in F$ be a point of a derived stack $F$. The \emph{tangent complex}
$T_xF$ is the derived mapping space
\[
  T_xF := \Map_{\mathrm{dStk}}(\Spec k[\epsilon]/(\epsilon^2), F)_x,
\]
considered as a complex in degrees $(-\infty,\dots,\infty)$.
\end{definition}

The cohomology groups of $T_xF$ encode:
\begin{itemize}
  \item $H^0(T_xF)$ classical tangent directions,
  \item $H^{-1}(T_xF)$ infinitesimal automorphisms,
  \item $H^1(T_xF)$ obstructions to extending deformations,
  \item higher degrees controlling higher obstructions and symmetries.
\end{itemize}

In particular, if $F$ is smooth (in derived sense), then $H^1(T_xF)$ may vanish,
ensuring unobstructed deformations; if $H^{-1}(T_xF)$ is nonzero, $x$ has
nontrivial automorphisms (gauge symmetries).

\Intuition
The tangent complex of $\Plenum$ will encode:
\begin{itemize}
  \item infinitesimal variations of $(\PhiField,\vField,\SField)$,
  \item gauge symmetries mixing $\PhiField$ and $\vField$,
  \item obstructions related to entropic singularities.
\end{itemize}

%---------------------------------------------------------------
\section{Cotangent Complex and Shifted Symplectic Structures}
%---------------------------------------------------------------

The cotangent complex $L_F$ of a derived stack $F$ is a fundamental invariant
generalizing the classical cotangent bundle. It plays a central role in the
definition of shifted symplectic structures.

\begin{definition}
The \emph{cotangent complex} $L_F$ is a perfect complex on $F$ (when $F$ is
derived Artin) whose dual $T_F := L_F^\vee$ generalizes the tangent bundle.
\end{definition}

Shifted symplectic structures are defined as certain closed bilinear forms on
$L_F$ (or $T_F$) of cohomological degree $k$. The case $k=-1$ is the setting of
the BV formalism and will be central to \rsvpabbr{}.

We discuss shifted symplectic structures in detail in
\cref{chap:shifted-symplectic}.

%---------------------------------------------------------------
\section{\texorpdfstring{$\infty$}{Infinity}--Groupoids and Fields}
%---------------------------------------------------------------

In ordinary geometry, a point in a moduli space of fields corresponds to a
single classical solution. In derived geometry, such points carry higher
homotopical structure: their automorphism groups form $\infty$--groupoids,
which encode gauge equivalences and higher symmetries.

Consequently, the moduli stack $\Plenum$ should be thought of as an object in
an $\infty$--category of derived stacks, with points representing field
configurations, morphisms representing gauge transformations, and higher
morphisms encoding redundancy among transformations themselves.

This viewpoint is essential to the BV formalism: ghosts, antifields, and the
master equation all arise naturally from the derived structure of $\Plenum$.

%---------------------------------------------------------------
\section{Summary and Outlook}
%---------------------------------------------------------------

This chapter introduced the basic machinery underlying derived geometry:
model categories, $\infty$--categories, $\infty$--topoi, and derived stacks,
with an emphasis on tangent and cotangent complexes and higher automorphisms.
These tools will be used in \cref{chap:shifted-symplectic,chap:aksz-bv} to
define and analyze shifted symplectic structures and the AKSZ formalism.

In \cref{chap:plenum-stack}, we will construct the lamphrodynamic moduli
stack $\Plenum$ explicitly as a derived Artin stack, then equip it with a
$(-1)$--shifted symplectic form in
\cref{chap:aksz-rsvp}. The present chapter thus serves as the homotopical
foundation for all subsequent developments.



% File: Volume-I-Foundations/chapters/ch03-shifted-symplectic.tex

\chapter{Shifted Symplectic Structures}
\label{chap:shifted-symplectic}

This chapter develops the basic notions of shifted symplectic geometry used to
formulate the moduli stack $\Plenum$ of lamphrodynamic configurations as a
$(-1)$--shifted symplectic derived Artin stack. We assume familiarity with the
cotangent complex of a derived stack and its dual tangent complex as introduced
in \cref{chap:inf-cat}. Throughout, we follow \cite{PTVV2013}, \cite{Calaque2015}
and related sources.

%--------------------------------------------------------------------
\section{Preliminaries: Classical Symplectic Geometry}
%--------------------------------------------------------------------

We begin with a brief review of classical symplectic geometry to fix notation
and highlight the generalization to derived settings.

\begin{definition}
A \emph{symplectic manifold} is a pair $(X,\omega)$ where $X$ is a smooth
manifold and $\omega \in \Omega^2(X)$ is a closed (i.e.\ $\dd\omega = 0$) and
nondegenerate $2$--form.
\end{definition}

A central example is $T^{\ast}Q$ for a smooth manifold $Q$, equipped with the
canonical $2$--form $\omega = \dd p_i \wedge \dd q^i$. Classical Hamiltonian
mechanics is framed on $(X,\omega)$ with Hamiltonian $H\in C^\infty(X)$.

In gauge theory and in field theory, however, the moduli space of classical
solutions is typically singular, infinite dimensional, and equipped with
higher automorphisms, so classical manifolds are no longer adequate. A
homotopical refinement is needed, leading to the notion of shifted symplectic
structures.

%--------------------------------------------------------------------
\section{Graded Geometry and Cohomological Degree}
%--------------------------------------------------------------------

In derived geometry, differential forms are replaced by cohomological complexes
and their degrees must be tracked carefully. A $k$--shifted symplectic form is,
informally, a closed $2$--form of cohomological degree $k$ on the derived
stack. Heuristically:
\[
  \omega \in H^{k}\big(\mathcal{A}^{2,cl}(F)\big),
\]
where $F$ is a derived Artin stack and $\mathcal{A}^{2,cl}(F)$ is the sheaf of
closed $2$--forms on $F$.

\begin{definition}
Let $F$ be a derived Artin stack. A \emph{$k$--shifted symplectic structure} on
$F$ is a closed $2$--form $\omega$ of cohomological degree $k$, represented by
a morphism of complexes
\[
  \omega: \wedge^2 L_F \longrightarrow \mathcal{O}_F[k],
\]
where $L_F$ is the cotangent complex of $F$, such that the induced pairing
$L_F^\vee \otimes L_F^\vee \to \mathcal{O}_F[k]$ is nondegenerate in the
derived sense.
\end{definition}

Here $L_F^\vee = T_F$ generalizes the tangent bundle. The condition that $L_F$
be perfect (locally of finite Tor dimension) ensures that $T_F$ behaves well as
a dual.

\Intuition
The shift $k$ controls the ``oddness'' or ``evenness'' of the symplectic form.
Classical symplectic structures correspond to $k=0$. The case $k=-1$
corresponds to odd symplectic geometry and plays the role of BV phase spaces.

%--------------------------------------------------------------------
\section{$(-1)$--Shifted Symplectic Structures and BV}
%--------------------------------------------------------------------

The case $k=-1$ is the setting of the Batalin–Vilkovisky (BV) formalism for
gauge theories. In this context, the derived stack $F$ encodes classical fields
and gauge symmetries; the shifted symplectic form $\omega$ is odd; and the BV
bracket
\[
  \{-,-\}: \mathcal{O}_F \otimes \mathcal{O}_F \to \mathcal{O}_F[-(-1)] = \mathcal{O}_F[1]
\]
is defined from $\omega$.

\begin{definition}
Given a $(-1)$--shifted symplectic structure $\omega$ on $F$, the
\emph{Batalin–Vilkovisky bracket} $\{-,-\}$ on functions (or, more precisely,
on appropriate cochains) is the unique odd Poisson bracket induced by $\omega$.
\end{definition}

The BV mechanism encodes gauge symmetries as degenerate directions of the
action, formally resolved by ghosts and higher homotopies. The \emph{classical
master equation}
\[
  \{S,S\} = 0
\]
then expresses invariance under these symmetries. Later, in
\cref{chap:aksz-bv,chap:aksz-rsvp}, we will show how the lamphrodynamic action
$S_{\rsvpabbr}$ arises from an AKSZ construction on the moduli stack
$\Plenum$ equipped with a $(-1)$--shifted symplectic structure.

%--------------------------------------------------------------------
\section{Closedness and Nondegeneracy in the Derived Sense}
%--------------------------------------------------------------------

Closedness of $\omega$ is expressed in terms of the differential and the
cohomological structure on $L_F$. Nondegeneracy means that the induced morphism
\[
  T_F \longrightarrow L_F[-k]
\]
is an isomorphism in the derived category. For $k=-1$, this becomes
\[
  T_F \longrightarrow L_F[1],
\]
which is automatically compatible with the BV grade.

\begin{remark}
For $k<0$, the shifted symplectic form is ``odd'' and induces higher order
compatibilities among tangent and cotangent complexes, matching precisely the
structure required for the BV bracket.
\end{remark}

Closedness means that $\omega$ lifts to a cocycle in a suitable truncated
de~Rham complex on $F$. We do not reproduce the full homotopical definition here
(see \cite{PTVV2013}), but emphasize that closedness is a global condition
involving sheaf descent and higher coherences.

%--------------------------------------------------------------------
\section{Examples and Fundamental Results}
%--------------------------------------------------------------------

Key results from Pantev–Toën–Vaquié–Vezzosi \cite{PTVV2013} include:

\begin{theorem}[PTVV]
Let $X$ be a smooth scheme and let $\Tstar{k}X$ be the $k$--shifted cotangent
stack. Then $\Tstar{k}X$ carries a canonical $k$--shifted symplectic structure.
\end{theorem}

This theorem provides the fundamental example of shifted symplectic structures
and plays a role analogous to the cotangent bundle in classical mechanics. In
derived geometry, one generalizes the phase space to $\Tstar{k}X$ for any
integer $k$, with $k=-1$ corresponding to BV theory.

Another key result concerns mapping stacks:

\begin{theorem}[PTVV]
Let $M$ be an oriented smooth manifold of dimension $d$, and let $F$ be a
derived Artin stack with an $n$--shifted symplectic structure. Then the
mapping stack $\Map(M,F)$ carries an $(n-d)$--shifted symplectic structure.
\end{theorem}

This result is central to the AKSZ construction, where $M$ is typically the
source manifold (e.g.\ the spacetime worldvolume) and $F$ is the target stack
encoding the fields. For \rsvpabbr{}, $M$ will often be spacetime or a
spacelike hypersurface, and $F$ the lamphrodynamic stack $\Plenum$.

\Intuition
Shifted symplecticity propagates from target to mapping stack, reducing the
shift by the dimension of the domain. For a $3$--dimensional worldvolume and
an $n$--shifted target, the resulting structure is $(n-3)$--shifted.

%--------------------------------------------------------------------
\section{Shifted Symplectic Forms from Field Content}
%--------------------------------------------------------------------

In \rsvpabbr{}, the $(-1)$--shifted symplectic structure arises from the
lamphrodynamic triplet $(\PhiField,\vField,\SField)$ and their coupling to
geometry. Roughly speaking, the moduli stack $\Plenum$ of field configurations
admits a cotangent complex generated by variations of the fields. Pairing these
variations through curvature and entropy flow produces a candidate closed
$3$--form, whose degree and symmetries align with $k=-1$.

In \cref{chap:plenum-stack}, we will:
\begin{enumerate}
  \item Construct $\Plenum$ as a derived Artin stack.
  \item Identify its cotangent complex $L_{\Plenum}$ explicitly in terms of
        variations of $(\PhiField,\vField,\SField)$.
  \item Construct a $(-1)$--shifted symplectic form $\omega_{\Plenum}$
        expressed locally as a derived $3$--form built from the lamphrodynamic
        fields.
  \item Verify that $\omega_{\Plenum}$ is closed and nondegenerate.
\end{enumerate}

This is a central technical achievement of Volume~I.

%--------------------------------------------------------------------
\section{Why \texorpdfstring{$k=-1$}{k = -1}?}
%--------------------------------------------------------------------

A natural question is: why specifically $k=-1$ for \rsvpabbr{}? Why not $k=0$
or $k=-2$?

\begin{itemize}
  \item The case $k=0$ would produce an \emph{even} symplectic form. This is
  appropriate for classical mechanics or certain topological sigma models, but
  not for gauge theories with ghost content.

  \item The BV formalism requires an odd symplectic form so that the BV bracket
  has the correct parity. This forces $k$ to be odd and negative.
\end{itemize}

Moreover, gauge symmetries and redundancy among lamphrodynamic fields indicate
that $\Plenum$ is not a naive moduli space but a derived stack with nontrivial
homotopy in negative degrees. The canonical choice compatible with BV is
$k=-1$, aligning with the degree of the ghost sector and with the reductions
arising in AKSZ.

\begin{remark}
In some exotic generalizations, higher negative shifts may appear (e.g.
topological field theories or higher spin models). For \rsvpabbr{}, the
entropic and lamphrodynamic structure strongly suggest $k=-1$ as the natural
choice.
\end{remark}

%--------------------------------------------------------------------
\section{Shifted Symplecticity and AKSZ}
%--------------------------------------------------------------------

The AKSZ construction, introduced in \cite{AKSZ}, produces solutions of the
classical master equation from a target derived stack equipped with a shifted
symplectic structure. The recipe is:
\begin{enumerate}
  \item Choose a source manifold $M$ (e.g.\ spacetime).
  \item Choose a target derived stack $F$ with an $n$--shifted symplectic
    structure.
  \item Form the mapping stack $\Map(M,F)$; by PTVV, this inherits an
    $(n-\dim M)$--shifted symplectic structure.
  \item For $n-\dim M = -1$, obtain an odd symplectic structure on
    $\Map(M,F)$, suitable for BV quantization.
\end{enumerate}

In \rsvpabbr{}, $F$ will ultimately be the lamphrodynamic target $\Plenum$ and
$M$ the spacetime manifold. The resulting $(-1)$--shifted symplectic structure
on $\Map(M,\Plenum)$ provides the geometric foundation for the BV action
$S_{\rsvpabbr}$.

%--------------------------------------------------------------------
\section{Summary and Forward Outlook}
%--------------------------------------------------------------------

Shifted symplectic geometry generalizes classical symplectic structures to a
homotopical context and provides the correct language for BV and AKSZ
quantization. The case $k=-1$ corresponds precisely to odd symplectic geometry
required for lamphrodynamic gauge theory.

In the next chapter, \cref{chap:aksz-bv}, we will review the AKSZ formalism in
detail, emphasizing how shifted symplectic structures on target stacks
determine BV actions, and how the constraint $n-\dim M=-1$ leads naturally to
lamphrodynamic dynamics on $\Plenum$.


\chapter{AKSZ Formalism and the Classical BV Construction}

\section{Introduction and Motivation}

This chapter introduces the Alexandrov--Kontsevich--Schwarz--Zaboronsky
(AKSZ) construction for classical field theories with
a $(k-1)$-shifted symplectic target, following
\cite{AKSZ97,CF01,CF07,PTVV13}.
The formalism generalizes the classical Batalin–Vilkovisky (BV) method
and provides a canonical route to constructing solutions of the
classical master equation. The approach is intrinsically geometric,
relying on mapping stacks, graded manifolds,
and homological vector fields of degree $+1$.

From the perspective of RSVP, the AKSZ framework is essential.
RSVP fields will be realized as maps
\[
\mathcal{F} \colon \mathsf{T}[1]M \longrightarrow \mathcal{X},
\]
where $\mathcal{X}$ is the derived moduli stack of
lamphrodynamic triples $(\Phi,\vec v,S)$ constructed in Chapter~6,
equipped with a canonical $(-1)$-shifted symplectic structure.
The AKSZ recipe then produces a BV action functional $S_{\text{BV}}$
satisfying $\{S_{\text{BV}},S_{\text{BV}}\}=0$, thereby yielding a
consistent classical theory in the BV sense.

This chapter develops the required formal background, establishes the
general AKSZ construction, and explains the classical BV complex.
The final section illustrates the formalism with the
example of three-dimensional Chern--Simons theory.

\medskip

\noindent{\bf Pedagogical note.}
The first half of this chapter is written from the perspective of a
physicist learning derived geometry; the second half assumes familiarity
with shifted symplectic structures, homological vector fields,
and the BV spectral sequence.


\section{Graded Manifolds and Homological Vector Fields}

Let $\mathcal{M}$ be a $\mathbb{Z}$-graded supermanifold (or more
generally, a derived stack locally modeled by such).
A \emph{homological vector field}
is a degree $+1$ vector field $Q$ on $\mathcal{M}$ satisfying
\[
Q^2 = \tfrac12[Q,Q]=0.
\]
The pair $(\mathcal{M},Q)$ is called a \emph{$Q$-manifold}.
The condition $Q^2=0$ implies that $Q$ determines a differential on
the algebra of functions $\mathcal{O}_{\mathcal{M}}$, and hence gives
rise to a cochain complex
$(\mathcal{O}_{\mathcal{M}},Q)$.

\begin{definition}
A \emph{$Q$-mapping stack} from $\mathsf{T}[1]M$ to $\mathcal{X}$
is a mapping stack
\[
\mathrm{Map}(\mathsf{T}[1]M,\mathcal{X})
\]
equipped with a homological vector field induced from the
vector fields on both source and target, in particular
the de Rham differential on $\mathsf{T}[1]M$.
\end{definition}

\begin{remark}
Physically, the degree shift encodes ghost number.
Coordinates of degree $0$ describe classical fields,
degree $-1$ antifields, degree $+1$ ghosts.
\end{remark}


\section{Shifted Symplectic Structures and the BV Bracket}

Let $(\mathcal{X},\omega)$ be a $k$-shifted symplectic derived stack.
Thus $\omega$ is a closed, non-degenerate $2$-form of cohomological
degree $k$.
When $k=-1$, the induced Poisson bracket has cohomological degree $+1$,
and is known as the BV bracket.

\begin{definition}
Given a $(-1)$-shifted symplectic structure $\omega$,
the \emph{BV bracket} $\{-,-\}$ on
$\mathcal{O}_{\mathcal{X}}$ is defined by
\[
\{f,g\} := \iota_{X_f}\iota_{X_g}\omega,
\]
where $X_f$ is the Hamiltonian vector field determined by $f$.
\end{definition}

\begin{proposition}
The BV bracket has the following properties:
\begin{enumerate}
\item $\{f,g\} = -(-1)^{(|f|+1)(|g|+1)}\{g,f\}$,
\item $\{f,gh\} = \{f,g\}h + (-1)^{(|f|+1)|g|}g\{f,h\}$,
\item $(-1)^{(|f|+1)(|h|+1)}\{f,\{g,h\}\} + \text{cyclic perms} = 0$.
\end{enumerate}
\end{proposition}

\begin{proof}[Proof (sketch)]
The identities follow from the graded symplectic structure and the
Cartan formula. A complete proof may be found in
\cite{Schwarz93,CattaneoFelder02,PTVV13}.
\end{proof}


\section{AKSZ Data and the Classical Master Equation}

Let $M$ be a smooth oriented manifold of dimension $d$,
and let $(\mathcal{X},\omega)$ be a $(d-1)$-shifted symplectic target
with a homological vector field $Q$ preserving $\omega$.
The AKSZ construction associates to this data a BV action on the
mapping stack
\[
\mathcal{F} := \mathrm{Map}(\mathsf{T}[1]M,\mathcal{X}),
\]
with an induced $(-1)$-shifted symplectic structure.
Concretely, the AKSZ action is given by
\begin{equation}
S_{\mathrm{AKSZ}} = \int_{\mathsf{T}[1]M}
\left( \langle \Phi^{\ast}\theta, d_M \Phi\rangle
    + \Phi^{\ast}\alpha \right),
\label{eq:AKSZ-action}
\end{equation}
where $\theta$ is a potential of $\omega$ and $\alpha$ is a Hamiltonian
function satisfying $Q=\{\alpha,-\}$.

\begin{theorem}[AKSZ classical master equation]
The functional $S_{\mathrm{AKSZ}}$ satisfies
\[
\{S_{\mathrm{AKSZ}},S_{\mathrm{AKSZ}}\}=0.
\]
\end{theorem}

\begin{proof}[Proof (sketch)]
The construction of $S_{\mathrm{AKSZ}}$ is arranged so that
the BV bracket encodes the graded symplectic form on the mapping space,
and the homological vector field encodes the differential on
$\mathsf{T}[1]M$.
The fact that $Q$ preserves $\omega$ and squares to zero implies
that $S_{\mathrm{AKSZ}}$ is a solution of the classical master equation.
For complete details see \cite{AKSZ97,CF01,CF07,PTVV13}.
\end{proof}

\section{Mapping Stacks and the Induced Shifted Symplectic Form}

Let $(\mathcal{X},\omega)$ be a $(d-1)$-shifted symplectic derived stack
in the sense of \cite{PTVV13}, and consider a compact oriented
$d$-manifold $M$. The AKSZ construction uses the internal hom
$\mathrm{Map}(\mathsf{T}[1]M,\mathcal{X})$, viewed as a derived mapping
stack in the $\infty$-topos of derived stacks. The tangent complex
is obtained by the standard formula
\[
T_{\Phi}\mathrm{Map}(\mathsf{T}[1]M,\mathcal{X})
\simeq \mathbf{R}\Gamma(\mathsf{T}[1]M, \Phi^{\ast}T_{\mathcal{X}}),
\]
which identifies tangent vectors at a field configuration $\Phi$
with sections of the pullback of the tangent complex of $\mathcal{X}$.

\begin{theorem}[PTVV]
\label{thm:PTVV-induced}
Let $(\mathcal{X},\omega)$ be $(d-1)$-shifted symplectic
and let $M$ be an oriented $d$-manifold.
Then the mapping stack $\mathrm{Map}(\mathsf{T}[1]M,\mathcal{X})$
carries a canonical $(-1)$-shifted symplectic structure
$\omega_{\mathrm{AKSZ}}$ induced by integration along $M$.
\end{theorem}

\begin{proof}[Proof (sketch)]
The degree shift arises from the Thom form of $M$ and the graded geometry
of $\mathsf{T}[1]M$. Integration reduces degree by $d$, and combined with
the $(d-1)$-shift of $\omega$, yields $-1$. The construction and closedness
are established in \cite{PTVV13}.
\end{proof}

\begin{remark}
This result is fundamental:
any $(d-1)$-shifted symplectic target yields a $(-1)$-shifted symplectic
BV theory on the mapping stack. In particular, the RSVP target constructed
in Chapter~6 will have $d=3$, hence a $(-1)$-shifted BV theory in
four spacetime dimensions.
\end{remark}


\section{BV Fields, Ghosts, and Antifields}

Given the AKSZ mapping stack
$\mathcal{F}=\mathrm{Map}(\mathsf{T}[1]M,\mathcal{X})$,
the derived geometry automatically organizes field components by ghost
number. More precisely, a degree-$k$ coordinate on $\mathcal{X}$ pulls
back to a degree-$(k+i)$ coordinate on $\mathcal{F}$, where $i$ labels
the internal grading of $\mathsf{T}[1]M$.
Classical fields occur at total degree $0$, ghosts at $+1$, antifields at
$-1$, and higher ghosts at degrees $>+1$.

\begin{definition}
The \emph{BV complex} associated to $\mathcal{F}$ is the cochain complex
$(\mathcal{O}_{\mathcal{F}},Q_{\mathrm{AKSZ}})$, where
$Q_{\mathrm{AKSZ}}$ is the homological vector field induced from both the
target differential and the de~Rham differential on the source.
\end{definition}

\begin{remark}
In classical physics language, we may write fields schematically as
\[
\{ \text{fields } \phi,\;\text{ghosts } c,\;\text{antifields }
\phi^{\ast},\;\text{anti-ghosts } c^{\ast},\ldots \},
\]
with degrees organized by the AKSZ grading.
\end{remark}


\section{The Classical BV Action and the Master Equation}

Let $(\mathcal{X},\omega)$ be $(d-1)$-shifted symplectic, and let
$\theta$ be a potential of $\omega$ (i.e.\ $\omega = d\theta$).
Assume $\mathcal{X}$ carries a Hamiltonian function $\alpha$ of degree
$d$ such that the associated Hamiltonian vector field $Q=\{\alpha,-\}$
coincides with the homological vector field on $\mathcal{X}$. Then
the AKSZ action \eqref{eq:AKSZ-action} satisfies the BV classical
master equation
\[
\{S_{\mathrm{AKSZ}},S_{\mathrm{AKSZ}}\} = 0.
\]
This is the precise sense in which AKSZ produces a BV theory.

\begin{proposition}
The BV bracket associated to $\omega_{\mathrm{AKSZ}}$ coincides with
the variational Schouten bracket on functionals over the mapping stack.
\end{proposition}

\begin{proof}[Proof (sketch)]
This is essentially the compatibility of the PTVV bracket with the internal
Hom structure and the universal property of the mapping stack. See
\cite{PTVV13,CF07}.
\end{proof}


\section{Example: Three-Dimensional Chern--Simons Theory}

We illustrate the general AKSZ recipe by recovering three-dimensional
Chern--Simons gauge theory, following \cite{AKSZ97,CF01,CF07}.
Let $M$ be a $3$-dimensional oriented manifold and let
$\mathfrak{g}$ be a finite-dimensional Lie algebra with invariant bilinear
form $\langle\cdot,\cdot\rangle$. Consider the graded vector space
\[
\mathfrak{g}[1],
\]
concentrated in degree $1$, i.e.\ coordinates have ghost number $1$.
Equip $\mathfrak{g}[1]$ with the homological vector field
\[
Q(x) = \tfrac12 [x,x],
\]
which is of degree $+1$ and squares to zero by the Jacobi identity.

\begin{proposition}
The space $\mathfrak{g}[1]$ carries a canonical
$2$-shifted symplectic form (hence $(d-1)=2$), given by the bilinear form
of degree~$2$ induced from $\langle\cdot,\cdot\rangle$.
\end{proposition}

\begin{proof}[Proof (sketch)]
The invariant bilinear form provides a non-degenerate pairing of appropriate
degree, and closedness follows from the Lie algebra structure. See
\cite{AKSZ97}.
\end{proof}

Applying the general AKSZ recipe with $d=3$ and target $\mathfrak{g}[1]$,
the mapping stack $\mathrm{Map}(\mathsf{T}[1]M,\mathfrak{g}[1])$ obtains
a $(-1)$-shifted symplectic structure.
Let $A$ denote the superfield obtained from the pullback of the graded
coordinate on $\mathfrak{g}[1]$; its degree-$0$ component is an ordinary
connection $1$-form on $M$, and higher-degree components encode ghosts
and antifields.

The AKSZ action takes the familiar form
\begin{equation}
S_{\mathrm{CS}} = \int_M \left\langle A \wedge dA
+ \tfrac13 A \wedge [A,A] \right\rangle,
\end{equation}
where products and brackets are understood in the graded sense.
One verifies that $\{S_{\mathrm{CS}},S_{\mathrm{CS}}\}=0$ using the
graded Jacobi identity and invariance of $\langle\cdot,\cdot\rangle$.

\begin{remark}
The usual Chern--Simons action is the restriction of $S_{\mathrm{CS}}$
to the classical degree-zero part of the AKSZ field.
All ghosts, antifields, and higher-degree components arise automatically
from the AKSZ mapping stack, rather than being introduced by hand.
\end{remark}

\medskip

\noindent\textbf{Forward link to RSVP.}
In Chapter~7 we apply the AKSZ recipe to the
$3$-shifted symplectic target $\mathrm{Plenum}$ of the lamphrodynamic
triples $(\Phi,\vec v,S)$, constructed explicitly in Chapter~6.
Because $\mathrm{Plenum}$ is $(3-1)=2$-shifted symplectic, the associated
mapping stack over $\mathsf{T}[1]M$ in four dimensions acquires a canonical
$(-1)$-shifted BV symplectic structure, producing the classical
RSVP BV action.




% ------------------------------------------------------------
% Part II: Lamphrodynamic Field Theory
% ------------------------------------------------------------
\part{Lamphrodynamic Field Theory}

\chapter{The Lamphrodynamic Plenum: Axiomatic Foundation}

\section{Motivation and Overview}

The purpose of this chapter is to establish a rigorous axiomatic framework
for the lamphrodynamic plenum, consisting of a triple
$(\Phi,\vec v,S)$ of scalar, vector, and entropy fields.
The physical interpretation of these fields has been discussed informally
in earlier writings; here we develop a formal foundation suitable for
derived geometry and subsequent quantization via the AKSZ construction.

The lamphrodynamic plenum is conceived as a universal medium
of information and energetic exchange.
The scalar potential $\Phi$ encodes lamphrodynamic potential energy,
the vector field $\vec v$ represents oriented transport flows, and
the entropy field $S$ accounts for dissipative and diffusive structure.
The fundamental equations couple these fields in a manner that respects
derived symplectic geometry, thermodynamic irreversibility, and
causal structure.

The classical field theory of $(\Phi,\vec v,S)$ will be defined by a
variational principle in Section~5.5, producing Euler--Lagrange
equations whose analytical properties will be investigated in
Chapters~9--11. Our goal in this chapter is to introduce the field
variables, state the axioms that govern their behavior, and derive
initial differential identities and structural constraints.


\section{Definition of the Lamphrodynamic Fields}

Let $(M,g)$ be a smooth oriented time-oriented Lorentzian manifold
of dimension $4$. The lamphrodynamic fields consist of:

\begin{enumerate}
\item A real-valued scalar field $\Phi \in C^{\infty}(M,\mathbb{R})$.
\item A spacetime vector field $\vec v \in \Gamma(TM)$.
\item A real-valued entropy scalar field $S \in C^{\infty}(M,\mathbb{R})$.
\end{enumerate}

We write collectively
\[
\mathcal{U} := (\Phi,\vec v,S).
\]

\begin{axiom}[Field regularity]
For the purposes of the classical theory developed in this volume,
we assume $\Phi$, $S$ are smooth and $\vec v$ is a smooth vector field
on $M$. Weak and distributional extensions will be introduced in
Chapter~11.
\end{axiom}

\begin{remark}
In certain applications we will regard $S$ as taking values in
$\mathbb{R}_{\ge 0}$, though for the mathematical theory it suffices to
treat $S$ as real-valued.
\end{remark}


\section{Lamphrodynamic Coupling Principles}

The physical behavior of the scalar, vector, and entropy fields derives
from the principle that lamphrodynamics is governed by coupled 
transport, diffusion, and potential flow. We impose the following axioms.

\begin{axiom}[Transport--Potential coupling]
The vector field $\vec v$ transports the scalar potential $\Phi$
according to the advection equation
\begin{equation}
\label{eq:advect}
\mathcal{L}_{\vec v}\Phi = -\kappa_{\Phi}\Delta_g \Phi + R_{\Phi},
\end{equation}
where $\kappa_{\Phi}\ge 0$ is a diffusion coefficient,
$\Delta_g$ is the Laplace--Beltrami operator, and $R_{\Phi}$
is a nonlinear lamphrodynamic coupling term defined in
Section~5.6.
\end{axiom}

\begin{axiom}[Entropy production]
The entropy field $S$ satisfies
\begin{equation}
\label{eq:entropy}
\mathcal{L}_{\vec v}S = \kappa_{S}\Delta_g S + \sigma(\Phi,\vec v,S),
\end{equation}
where $\sigma(\Phi,\vec v,S)\ge 0$ is the entropy production rate.
\end{axiom}

\begin{axiom}[Velocity response]
The transport field $\vec v$ responds to both entropy and potential
gradients by
\begin{equation}
\label{eq:velocity}
\nabla_{\vec v}\vec v = -\nabla\Phi + \beta\nabla S + R_{v},
\end{equation}
where $\beta\in\mathbb{R}$ couples entropy gradients to lamphrodynamic
flows and $R_{v}$ is described in Section~5.7.
\end{axiom}

These axioms are formal for now; precise variational origins will be
developed in Section~5.5. The choice of signs corresponds to the
interpretation of $-\nabla\Phi$ as attractive potential flow and
$\beta\nabla S$ as repulsive entropic response.


\section{Lamphrodynamic Conservation and Balance Laws}

We now derive several differential identities and balance laws implied
by the axioms of the previous section. For clarity, we work in local
coordinates with the Levi--Civita connection of the Lorentzian metric $g$.

\begin{proposition}[Entropy balance]
Assume equation \eqref{eq:entropy}. Then
\begin{equation}
\partial_t \int_{\Sigma_t} S \, d\mu_t
= \int_{\Sigma_t} \left( \kappa_{S}\Delta_g S
  + \sigma(\Phi,\vec v,S) \right) d\mu_t
\end{equation}
for any spacelike hypersurface $\Sigma_t$ with induced measure $d\mu_t$.
\end{proposition}

\begin{proof}[Proof]
Apply the Reynolds transport theorem to \eqref{eq:entropy} and use
$\mathcal{L}_{\vec v}S = \vec v(S)$.
\end{proof}

Similarly, we have a potential energy balance:

\begin{proposition}[Potential balance]
If $\Phi$ satisfies \eqref{eq:advect}, then
\begin{equation}
\partial_t \int_{\Sigma_t} \Phi \, d\mu_t
= \int_{\Sigma_t}
\left( -\kappa_{\Phi}\Delta_g \Phi + R_{\Phi} \right) d\mu_t.
\end{equation}
\end{proposition}

The vector field $\vec v$ obeys a momentum-like balance obtained by
integrating \eqref{eq:velocity} against $\vec v$ and using the identity
$\frac12\nabla_{\vec v}\|\vec v\|^2 = \langle\nabla_{\vec v}\vec v,\vec v\rangle$.

\begin{proposition}[Lamphrodynamic momentum balance]
Assuming \eqref{eq:velocity},
\begin{equation}
\frac12 \mathcal{L}_{\vec v}\|\vec v\|^2 =
-\langle\nabla\Phi,\vec v\rangle
+ \beta\langle\nabla S,\vec v\rangle
+ \langle R_{v},\vec v\rangle.
\end{equation}
\end{proposition}

\begin{remark}
The quantity $\langle\nabla S,\vec v\rangle$ plays the role of
entropic tension, corresponding physically to emergent repulsion
induced by entropy gradients. This term will be crucial in the
cosmological applications of Volume~II.
\end{remark}

\section{Variational Principle and Euler--Lagrange Equations}
\label{sec:variational}

We now give the variational origin of the lamphrodynamic equations.
Let $\mathcal{L}$ be a Lagrangian density on $(M,g)$ of the form
\begin{equation}
\label{eq:Lagrangian}
\mathcal{L}(\Phi,\vec v,S;g)
= \tfrac12\|\vec v\|^2
+ U(\Phi) + W(S) + \alpha\langle\nabla S,\vec v\rangle
+ \beta\langle\nabla\Phi,\vec v\rangle,
\end{equation}
where $U$ and $W$ are smooth potential functions and
$\alpha,\beta\in\mathbb{R}$ are coupling constants. Let
\[
\mathcal{S}[\Phi,\vec v,S]
:= \int_M \mathcal{L}\, dV_g
\]
be the associated action functional.
Variation against $\Phi$, $S$, and $\vec v$ produces
Euler--Lagrange equations which we now compute formally.

\vspace{0.5em}
\noindent{\bf Variation with respect to $\Phi$.}
Integrating by parts yields
\[
\delta_{\Phi}\mathcal{S}
= \int_M \left(
U'(\Phi) + \beta \operatorname{div}(\vec v)
\right)\delta\Phi\, dV_g,
\]
and hence the Euler--Lagrange equation is
\begin{equation}
\label{eq:EL-Phi}
U'(\Phi) + \beta\operatorname{div}(\vec v) = 0.
\end{equation}

\vspace{0.5em}
\noindent{\bf Variation with respect to $S$.}
Similarly,
\[
\delta_{S}\mathcal{S}
= \int_M \left(
W'(S) + \alpha \operatorname{div}(\vec v)
\right)\delta S\, dV_g,
\]
producing
\begin{equation}
\label{eq:EL-S}
W'(S) + \alpha\operatorname{div}(\vec v)=0.
\end{equation}

\vspace{0.5em}
\noindent{\bf Variation with respect to $\vec v$.}
A computation using the Levi--Civita connection gives
\[
\delta_{\vec v}\mathcal{S}
= \int_M
\left(
\vec v + \alpha\nabla S + \beta\nabla\Phi
\right)\cdot \delta\vec v\, dV_g,
\]
leading to
\begin{equation}
\label{eq:EL-v}
\vec v + \alpha\nabla S + \beta\nabla\Phi = 0.
\end{equation}

\begin{remark}
Equations \eqref{eq:EL-Phi}, \eqref{eq:EL-S}, and \eqref{eq:EL-v}
represent the stationary points of the lamphrodynamic action.
Their consistency with the axioms of Section~5.3 requires suitable
choices of $U,W,\alpha,\beta$; Section~5.6 analyzes this structure.
\end{remark}


\section{Nonlinear Couplings and Constraint Relations}
\label{sec:couplings}

We now discuss the nonlinear terms $R_{\Phi}$ and $R_{v}$ introduced
formally in Section~5.3. These terms encode constraint relations among
$\Phi$, $\vec v$, and $S$ that are not captured by the minimal action
\eqref{eq:Lagrangian}. In physical terms, $R_{\Phi}$ and $R_{v}$
represent lamphrodynamic self-interaction and dissipation beyond the
lowest-order coupling.

\begin{axiom}[Lamphrodynamic constraint]
There exist smooth functions $\mathcal{C}_\Phi(\Phi,\vec v,S)$ and
$\mathcal{C}_v(\Phi,\vec v,S)$ such that
\begin{equation}
\label{eq:R-Phi}
R_{\Phi} = \mathcal{C}_\Phi(\Phi,\vec v,S),\qquad
R_{v} = \mathcal{C}_{v}(\Phi,\vec v,S),
\end{equation}
subject to the compatibility condition
\begin{equation}
\label{eq:compatibility}
\operatorname{div}(R_{v})
+ \langle\nabla S,R_{v}\rangle
+ \langle\nabla\Phi,R_{v}\rangle
= \Delta_g R_{\Phi}.
\end{equation}
\end{axiom}

\begin{remark}
Condition \eqref{eq:compatibility} guarantees that
Lamphrodynamic balance laws are preserved at nonlinear order,
and plays a role analogous to constitutive relations in hydrodynamics.
\end{remark}


\section{Symmetries and Noether Currents}
\label{sec:symmetry-noether}

We discuss the symmetries of the lamphrodynamic plenum at the level
of the classical action \eqref{eq:Lagrangian}.

\begin{definition}
A local symmetry of the lamphrodynamic action is a variation
$(\delta\Phi,\delta\vec v,\delta S)$ such that
$\delta\mathcal{S}=0$ for all field configurations satisfying
appropriate boundary conditions.
\end{definition}

\begin{proposition}[Shift symmetry of $\Phi$]
If $U$ depends only on $\nabla\Phi$ (e.g.\ pure gradient theory),
then the transformation $\Phi\mapsto \Phi+\epsilon$ is a symmetry
of $\mathcal{S}$, yielding a conserved current
\[
J^{\mu}_{\Phi} = \frac{\partial\mathcal{L}}{\partial(\nabla_{\mu}\Phi)}.
\]
\end{proposition}

\begin{remark}
In cosmological settings, $U$ will typically depend on $\Phi$ itself,
breaking shift symmetry and producing effective potential-induced flow.
\end{remark}

Likewise, if $W$ depends only on $\nabla S$, then $S\mapsto S+\epsilon$
is a symmetry; otherwise, entropy production breaks $S$-shift symmetry.
The behavior of these symmetries in the presence of $R_{\Phi}$ and
$R_v$ is subtle and will be discussed further in Volume~II.


\section{Conservation Laws and Energy Conditions}

We conclude this chapter with an analysis of energy conditions for
lamphrodynamic flows.
Let
\[
\mathcal{E} := \tfrac12\|\vec v\|^2 + U(\Phi) + W(S)
\]
denote the lamphrodynamic energy density.
Assuming equations \eqref{eq:EL-Phi}--\eqref{eq:EL-v}, one computes

\begin{proposition}[Lamphrodynamic energy balance]
\begin{equation}
\mathcal{L}_{\vec v}\mathcal{E}
= - \alpha\|\nabla S\|^2 - \beta\|\nabla\Phi\|^2
+ \langle R_{v},\vec v\rangle + R_{\Phi}.
\end{equation}
\end{proposition}

\begin{proof}[Proof]
Differentiating $\mathcal{E}$ along the flow of $\vec v$ and substituting
\eqref{eq:EL-Phi}--\eqref{eq:EL-v} yields the stated identity.
\end{proof}

\begin{remark}
If $\alpha,\beta>0$ and $R_{\Phi}=R_{v}=0$, then
$\mathcal{L}_{\vec v}\mathcal{E}\le 0$, indicating lamphrodynamic
dissipation. The nonlinear regime will be addressed in Chapter~10 via
a priori energy estimates.
\end{remark}

\medskip

\noindent\textbf{Summary and Preview.}
This chapter has introduced the lamphrodynamic fields, stated axioms for
their coupling, developed a variational formulation, and identified
initial balance laws and symmetries. In Chapter~6 we construct the
derived moduli stack of lamphrodynamic configurations and show that it
admits a canonical $(3-1)=2$-shifted symplectic structure,
crucial for the AKSZ quantization developed in Chapter~7.


\chapter{Derived Stack Construction for RSVP}

\section{Motivation and Overview}

In Chapter~5 the lamphrodynamic fields $(\Phi,\vec v,S)$ were introduced
as classical smooth fields on a Lorentzian spacetime $(M,g)$.
To quantize the theory using the AKSZ formalism,
we require a \emph{derived} moduli stack whose classical points correspond
to such triples, so that the mapping stack
\[
\mathrm{Map}(\mathsf{T}[1]M,\mathrm{Plenum})
\]
inherits a $(-1)$-shifted symplectic structure.
This chapter constructs the derived moduli stack $\mathrm{Plenum}$,
proves that it is a derived Artin stack locally of finite presentation,
describes its tangent complex and obstruction theory,
and prepares the ground for the shifted symplectic structure in
Chapter~7.

We rely on foundational results of Lurie \cite{Lurie-DAG} and
To\"en--Vezzosi \cite{ToenVezzosi-HAGII} on derived algebraic geometry,
and on Pantev--To\"en--Vaqui\'e--Vezzosi \cite{PTVV13} for the
construction of shifted symplectic forms on derived mapping stacks.


\section{The Derived Moduli Functor}

Fix a smooth oriented $4$-manifold $M$ with Lorentzian metric $g$.
For a derived affine scheme $\operatorname{Spec}A$,
define the space of lamphrodynamic fields over $A$ by
\[
\mathcal{U}(A)
:= \left\{
(\Phi,\vec v,S) \mid 
\Phi\in\Gamma(M_A,\mathcal{O}_{M_A}),
\;
\vec v\in\Gamma(M_A,T_{M_A}),
\;
S\in\Gamma(M_A,\mathcal{O}_{M_A})
\right\},
\]
where $M_A := M\times\operatorname{Spec}A$ and
$T_{M_A}$ is the relative tangent complex.
We regard $\mathcal{U}$ as a functor
\[
\mathcal{U} \colon \mathbf{dAff}^{\mathrm{op}}
\longrightarrow \mathbf{sSet}.
\]

\begin{definition}
The \emph{derived lamphrodynamic moduli stack} is the derived
stackification of $\mathcal{U}$, denoted
\[
\mathrm{Plenum} := \mathbf{L}\mathcal{U}.
\]
\end{definition}

Classical field configurations correspond to $A=\mathbb{R}$,
so that $\mathrm{Plenum}(\mathbb{R})$ recovers the classical triples
introduced in Chapter~5.


\section{Representability and Derived Artin Property}

We now establish that $\mathrm{Plenum}$ is a derived Artin stack.
Let us outline the representability argument; complete details follow
the general pattern of mapping stacks of sections of a vector bundle,
together with the derived enhancement required by the lamphrodynamic
constraint.

\begin{theorem}
\label{thm:Plenum-Artin}
The derived lamphrodynamic moduli stack $\mathrm{Plenum}$ is
a derived Artin stack locally of finite presentation over $\mathbb{R}$.
\end{theorem}

\begin{proof}[Proof (sketch)]
The moduli problem is locally given by sections $\Phi,S$ of trivial
line bundles and a section $\vec v$ of the tangent bundle.
Such mapping stacks are representable by derived Artin stacks
(\cite{ToenVezzosi-HAGII}, \S7).
The nonlinear constraints \eqref{eq:R-Phi}--\eqref{eq:compatibility}
are encoded as derived intersections of smooth Artin stacks,
preserving Artin representability. For details see
\cite{Lurie-DAG,ToenVezzosi-HAGII}.
\end{proof}

\begin{remark}
The lamphrodynamic constraint creates a derived structure:
classically one would impose $\mathcal{C}_\Phi=\mathcal{C}_v=0$,
but this would produce singularities in the moduli space.
Derived geometry resolves these singularities by enriching the
intersection with higher homotopical data.
\end{remark}


\section{Local Models and the Tangent Complex}

Given $\mathcal{U}$ as above,
the tangent complex at a point
\[
(\Phi_0,\vec v_0,S_0)\in \mathrm{Plenum}
\]
is given by the derived global sections
\begin{align}
T_{(\Phi_0,\vec v_0,S_0)}\mathrm{Plenum}
\simeq
\mathbf{R}\Gamma\bigl(
M,
\mathcal{O}_M \oplus T_M \oplus \mathcal{O}_M
\bigr).
\end{align}
When nonlinear constraints are included, the tangent complex becomes the
fiber of
\[
\mathbf{R}\Gamma(
M,
\mathcal{O}_M \oplus T_M \oplus \mathcal{O}_M
)
\longrightarrow
\mathbf{R}\Gamma(M, \mathcal{N}),
\]
where $\mathcal{N}$ is the derived normal bundle determined by the
Jacobian of the constraint functions
$\mathcal{C}_\Phi,\mathcal{C}_v$.

\begin{proposition}
The tangent complex of $\mathrm{Plenum}$ has perfect amplitude
contained in $[-1,1]$.
\end{proposition}

\begin{proof}[Proof (sketch)]
The tangent complex of sections of a perfect complex has perfect
amplitude given by the amplitude of the underlying complex of sheaves.
The derived constraint adds a fiber whose amplitude is bounded by $1$.
Standard arguments apply \cite{Lurie-DAG}.
\end{proof}

\begin{remark}
The amplitude $[-1,1]$ is crucial: it ensures $\mathrm{Plenum}$ has
a well-defined shifted symplectic structure of degree $2$; see Chapter~7.
\end{remark}

\section{Obstruction Theory and Derived Intersections}

We now analyze the derived enhancement induced by the nonlinear
lamphrodynamic constraints
\eqref{eq:R-Phi}--\eqref{eq:compatibility}.
Following \cite{ToenVezzosi-HAGII}, derived algebraic geometry replaces
classical intersections of moduli problems with homotopy fiber products.
Let $\mathcal{M}_{\Phi,S}$ denote the moduli stack of pairs of functions
$(\Phi,S)\in\Gamma(M,\mathcal{O}_M)^2$, and let $\mathcal{M}_{v}$ denote
the moduli stack of $\Gamma(M,T_M)$-valued sections.
The lamphrodynamic constraints define classical morphisms
\[
\mathcal{M}_{\Phi,S}\times \mathcal{M}_{v}
\longrightarrow \mathcal{N}_{\Phi,v}
\]
whose zero locus would be singular in general.

\begin{definition}
The \emph{derived lamphrodynamic intersection} is the homotopy fiber
product
\[
\mathrm{Plenum}
\simeq
(\mathcal{M}_{\Phi,S}\times \mathcal{M}_{v})
\times^{\mathbf{R}}_{\mathcal{N}_{\Phi,v}} \ast.
\]
\end{definition}

\begin{remark}
The derived homotopy fiber product retains information lost by
classical intersection, encoding obstructions in higher homotopical
degrees. This derived structure is responsible for the shifted
symplectic geometry we exploit in Chapter~7.
\end{remark}


\section{Derived Normal Bundle and Obstruction Classes}

Let $\mathcal{N}_{\Phi,v}$ denote the normal bundle determined by the
Jacobian of the lamphrodynamic constraint functions.
The derived normal bundle is obtained by taking the homotopy fiber
of the Jacobian mapping between perfect complexes, producing a derived
vector bundle in degrees $[-1,0]$.
The obstruction theory for $\mathrm{Plenum}$ is hence controlled by the
derived cohomology of $\mathcal{N}_{\Phi,v}$, localized on $M$.

\begin{theorem}
\label{thm:obstruction-Plenum}
The derived moduli stack $\mathrm{Plenum}$ admits a perfect obstruction
theory in the sense of \cite{BF97}, induced by the derived normal bundle
of the lamphrodynamic constraint.
\end{theorem}

\begin{proof}[Proof (sketch)]
The lamphrodynamic constraint defines a morphism of derived Artin stacks.
The homotopy fiber of the Jacobian defines a perfect obstruction theory,
using the general representability results of
\cite{BF97,ToenVezzosi-HAGII}.
\end{proof}

\begin{remark}
Physically, the obstruction theory describes infinitesimal lamphrodynamic
defects: perturbations of $(\Phi,\vec v,S)$ that are obstructed by the
constraint. In Volume~II these play a role near entropic singularities.
\end{remark}


\section{Existence of a $2$-Shifted Symplectic Structure}

We now show that $\mathrm{Plenum}$ admits a natural $2$-shifted
symplectic structure. The general existence theorem in \cite{PTVV13}
applies to derived intersections of moduli of sections of perfect
complexes when certain orientation data are available. In the present
setting, the Lorentzian manifold $(M,g)$ provides the necessary
orientation and volume form, and the kinetic and coupling terms in the
action provide a natural closed $3$-form.

Let $\Theta$ be the lamphrodynamic $3$-form defined locally by
\begin{equation}
\label{eq:local3form}
\Theta :=
\Phi\, dS\wedge \star d\Phi
+ \beta\, \langle\nabla\Phi,\nabla S\rangle\, dV_g
+ \gamma\, \langle\nabla S,\vec v\rangle\, d(\iota_{\vec v}dV_g),
\end{equation}
where $\star$ is the Hodge operator associated to $g$ and
$\gamma$ is a normalizing constant determined below.

\begin{proposition}
The form $\Theta$ extends to a closed $3$-form on the derived moduli
stack $\mathrm{Plenum}$, defining a $2$-shifted presymplectic structure.
\end{proposition}

\begin{proof}[Proof (sketch)]
Each term in \eqref{eq:local3form} is closed at the classical level
modulo the lamphrodynamic constraint. The derived extension is obtained
by pulling back to the derived intersection and verifying closedness in
the derived de~Rham complex, following \cite{PTVV13}. Full details will
be provided in Appendix~E.
\end{proof}


\section{Local Charts and Explicit Expression for $\omega$}

We now identify the $2$-shifted symplectic form $\omega$ as the derived
differential of $\Theta$:
\[
\omega := d_{\mathrm{dR}}\Theta.
\]
Classically, this yields
\begin{align}
\omega
&=
d\Phi\wedge dS\wedge \star d\Phi
+ \Phi\, d(dS\wedge \star d\Phi)
\\
& \qquad
+ \beta\, d\bigl(
\langle\nabla\Phi,\nabla S\rangle\, dV_g\bigr)
+ \gamma\, d\bigl(
\langle\nabla S,\vec v\rangle\, d(\iota_{\vec v}dV_g)
\bigr).
\end{align}
The derived extension incorporates infinitesimal deformations and ghosts
of each lamphrodynamic field. We omit the full expression here, as it
requires the total derived de~Rham complex on an Artin stack; see
Appendix~E.

\begin{theorem}[Shifted symplectic structure of $\mathrm{Plenum}$]
\label{thm:shifted}
The closed derived $3$-form $\Theta$ induces a
$2$-shifted symplectic structure $\omega$ on $\mathrm{Plenum}$.
\end{theorem}

\begin{proof}[Proof (sketch)]
Closedness is by construction. Non-degeneracy follows from the perfect
amplitude $[-1,1]$ of the tangent complex and the existence of
the lamphrodynamic constraint which produces a derived intersection of
section spaces satisfying the orientation axioms of \cite{PTVV13}.
See also \cite{ToenVezzosi-HAGII}.
\end{proof}

\begin{remark}[Physical origin of the shift]
Since the physical spacetime dimension is $4$,
integration of the derived $3$-form over the source $M$ produces an
effective shift by $3-4=-1$ in the BV symplectic degree, giving a
$(-1)$-shifted BV structure in Chapter~7. Hence the physical degree $k$
for RSVP quantization is forced to be $2$ at the classical level,
consistent with the AKSZ construction for $4$-dimensional theories.
\end{remark}

\section{Main Structural Theorem for the Lamphrodynamic Plenum}

We now unify the constructions of the preceding sections into a single
structural result, collecting assumptions and conclusions in one place.

\begin{theorem}[Main Structural Theorem]
\label{thm:Main-Plenum}
Let $M$ be a smooth oriented time-oriented $4$-dimensional Lorentzian
manifold and let $(\Phi,\vec v,S)$ be the lamphrodynamic fields with
nonlinear constraints \eqref{eq:R-Phi}--\eqref{eq:compatibility}.
Then the following properties hold:
\begin{enumerate}
\item The moduli problem of lamphrodynamic triples extends to a derived
functor $\mathcal{U}\colon\mathbf{dAff}^{\mathrm{op}}\to\mathbf{sSet}$
whose derived stackification $\mathrm{Plenum}$ is a derived Artin stack
locally of finite presentation.

\item The nonlinear constraints are resolved by a perfect obstruction
theory induced from the derived normal bundle of the constraint, and
$\mathrm{Plenum}$ has a tangent complex of perfect amplitude $[-1,1]$.

\item There exists a canonical closed derived $3$-form $\Theta$ whose
derived de~Rham differential $\omega := d_{\mathrm{dR}}\Theta$
defines a $2$-shifted symplectic structure on $\mathrm{Plenum}$.

\item The shifted symplectic structure is functorial with respect to
pullback along smooth morphisms of Lorentzian spacetimes,
and admits local presentations on Artin atlases of $\mathrm{Plenum}$
compatible with the lamphrodynamic constraint.
\end{enumerate}
\end{theorem}

\begin{proof}[Proof (sketch)]
Part (1) follows from representability of mapping stacks of sections of
perfect complexes and stability of Artin representability under derived
homotopy fiber products \cite{ToenVezzosi-HAGII}.
Part (2) follows from the construction of the derived normal bundle and
Behrend--Fantechi obstruction theories \cite{BF97}.
Part (3) is an instance of the PTVV construction of shifted symplectic
forms on derived mapping stacks with perfect obstruction theory
\cite{PTVV13}.
Part (4) follows from standard functoriality properties of shifted
symplectic geometry \cite{Calaque16}.
\end{proof}


\section{Physical Interpretation}

Although the shift by $2$ appears purely technical, it reflects the fact
that RSVP is a $4$-dimensional physical theory whose classical degrees
of freedom arise from a 3-form $\Theta$. The descent from $k=2$ to $k-4=-2$
under integration over the source would ordinarily yield a $(-2)$-shift,
but the AKSZ construction uses $\mathsf{T}[1]M$ in place of $M$,
producing the canonical $(-1)$-shift required for the BV bracket.
Thus the appearance of degree $2$ is a physical consequence of the
lamphrodynamic structure in four dimensions.

\begin{remark}
One may view the lamphrodynamic plenum as a universal ambient space of
physical fields in which gravitational and entropic interactions are
encoded as derived symplectic geometry. In this sense, $(\Phi,\vec v,S)$
constitute a pre-geometric substrate from which spacetime causality
emerges (developed further in Volume~II).
\end{remark}


\section{Preview of the AKSZ--RSVP Action}

In Chapter~7 we construct the AKSZ action associated to the
$2$-shifted symplectic derived stack $\mathrm{Plenum}$.
The result will be a $(-1)$-shifted symplectic BV theory on the
mapping stack $\mathrm{Map}(\mathsf{T}[1]M,\mathrm{Plenum})$,
providing a geometric and functorial route to quantization.
The Euler--Lagrange equations of the AKSZ action reproduce the
lamphrodynamic field equations of Chapter~5, up to additional ghost and
antifield terms required for BV consistency.

\begin{remark}
In the AKSZ formalism, no gauge-fixing or ad-hoc ghosts are introduced:
the ghost and antifield sectors appear automatically, determined by the
shifted symplectic geometry. This is essential for a physically
well-defined quantization of RSVP.
\end{remark}

\chapter{AKSZ--RSVP Action and Classical Master Equation}

\section{Overview and Strategy}

The goal of this chapter is to apply the AKSZ construction of
Chapter~4 to the derived lamphrodynamic plenum constructed in
Chapter~6. The result is a $(-1)$-shifted symplectic BV theory whose
classical Euler--Lagrange equations agree with the lamphrodynamic
equations of Chapter~5. This constitutes the precise mathematical
definition of the RSVP field theory at the classical level.

Let $(M,g)$ be a smooth oriented time-oriented $4$-dimensional Lorentzian
manifold with associated shifted tangent bundle $\mathsf{T}[1]M$.
Define the AKSZ field space
\[
\mathcal{F} := \mathrm{Map}(\mathsf{T}[1]M,\mathrm{Plenum}),
\]
a derived mapping stack equipped with the canonical $(-1)$-shifted
symplectic form of Chapter~4.
Fields in $\mathcal{F}$ consist of superfields whose degree-zero
components are the lamphrodynamic fields $(\Phi,\vec v,S)$ and whose
higher-degree components constitute ghosts, antifields, and higher
ghosts.


\section{AKSZ Data for RSVP}

The AKSZ construction requires:

\begin{enumerate}
\item a source $(\mathsf{T}[1]M,d_M)$ with homological vector field
$d_M$, the de~Rham differential on $M$ shifted by $1$;
\item a target $(\mathrm{Plenum},\omega)$ with $2$-shifted symplectic
form $\omega$ and homological vector field $Q_{\mathrm{Plenum}}$
preserving $\omega$;
\item a Hamiltonian function $\alpha$ on $\mathrm{Plenum}$ satisfying
$Q_{\mathrm{Plenum}}=\{\alpha,-\}$.
\end{enumerate}

Under these assumptions, the AKSZ recipe produces the BV action
\[
S_{\mathrm{AKSZ}} = \int_{\mathsf{T}[1]M}
\langle \Phi^{\ast}\theta, d_M\Phi\rangle + \Phi^{\ast}\alpha,
\]
where $\theta$ is a potential for the shifted symplectic form on
$\mathrm{Plenum}$.
Explicit formulas for $\theta$ will be given below (Section~\ref{sec:theta}).


\section{Derived 3-Form and BV Potential}

Recall from Chapter~6 the derived $3$-form $\Theta$ on $\mathrm{Plenum}$
given locally by
\[
\Theta =
\Phi\, dS\wedge \star d\Phi
+ \beta\, \langle\nabla\Phi,\nabla S\rangle\, dV_g
+ \gamma\, \langle\nabla S,\vec v\rangle\, d(\iota_{\vec v}dV_g).
\]
Let $\theta$ be any degree-$2$ potential for $\omega=d_{\mathrm{dR}}\Theta$
in the derived de~Rham complex of $\mathrm{Plenum}$. Such a potential
exists locally and may be constructed globally by descent, using the
Artin atlases of $\mathrm{Plenum}$.

\begin{remark}
In practice, $\theta$ may be chosen to be the $2$-form part of $\Theta$,
omitting terms that vanish under pullback to the mapping stack.
However, the existence of a global potential is essential for defining
the AKSZ action.
\end{remark}


\section{The RSVP AKSZ Action}
\label{sec:theta}

The AKSZ action is now defined by
\begin{equation}
\label{eq:RSVP-AKSZ}
S_{\mathrm{RSVP}}
:= \int_{\mathsf{T}[1]M}
\left(
\langle \Phi^{\ast}\theta, d_M\Phi\rangle
+ \Phi^{\ast}\alpha
\right).
\end{equation}
In local coordinates, writing $\Phi^{\ast}=(\Phi,\vec v,S)$ for the
pullback of the universal fields, we may express the integrand as
\begin{align}
\mathcal{L}_{\mathrm{AKSZ}}
&=
\langle d_M\Phi, \theta_{\Phi}\rangle
+ \langle d_M \vec v, \theta_{v}\rangle
+ \langle d_M S, \theta_{S}\rangle
+ \alpha(\Phi,\vec v,S),
\end{align}
where $\theta_{\Phi},\theta_{v},\theta_{S}$ are the components of the BV
potential paired with the corresponding superfield derivatives.

\begin{proposition}
The functional $S_{\mathrm{RSVP}}$ is well defined on the derived
mapping stack $\mathcal{F}$ and has total ghost number zero.
\end{proposition}

\begin{proof}[Proof (sketch)]
Total degree zero follows from degree counting: the degree of
$d_M$ is $+1$, the degree of $\theta$ is $2$, and the integration over
$\mathsf{T}[1]M$ reduces degree by $3$, yielding degree $0$ overall.
See \cite{AKSZ97,CF07}.
\end{proof}

\section{Classical Master Equation}

Let $\{-,-\}_{\mathrm{BV}}$ denote the BV antibracket induced by the
$(-1)$–shifted symplectic structure on $\mathcal{F}
= \mathrm{Map}(\mathsf{T}[1]M,\mathrm{Plenum})$. The classical master
equation is
\[
\{S_{\mathrm{RSVP}},S_{\mathrm{RSVP}}\}_{\mathrm{BV}} = 0.
\]

\begin{theorem}[Classical Master Equation]
\label{thm:CME-RSVP}
The RSVP AKSZ action \eqref{eq:RSVP-AKSZ} satisfies the classical master
equation
\[
\{S_{\mathrm{RSVP}},S_{\mathrm{RSVP}}\}_{\mathrm{BV}} = 0,
\]
and hence defines a classical BV theory.
\end{theorem}

\begin{proof}[Proof (Sketch)]
This follows from Theorem~\ref{thm:Main-Plenum}
and the general AKSZ theorem of
Alexandrov–Kontsevich–Schwarz–Zaboronsky:
whenever the target carries a $k$–shifted symplectic structure with
Hamiltonian $Q_{\mathrm{Plenum}}=\{\alpha,-\}$,
the induced action on
$\mathrm{Map}(\mathsf{T}[1]M,\mathrm{Plenum})$
satisfies the classical master equation and
induces the correct BV differential $Q=\{S,-\}$.
\end{proof}

\begin{remark}
Note that no gauge-fixing is used. The ghosts and antifields arise
canonically from the graded geometry of $\mathrm{Map}(\mathsf{T}[1]M,
\mathrm{Plenum})$ and the shifted symplectic structure.
\end{remark}


\section{BV Operator and Field Content}
\label{sec:BV-operator}

The BV differential is defined by
\[
Q_{\mathrm{BV}}(F) :=
\{ S_{\mathrm{RSVP}}, F\}_{\mathrm{BV}}, \qquad F\in
\mathcal{O}(\mathcal{F}).
\]
The field content consists of the superfields
\[
\mathbf{\Phi} =
\Phi + \Phi^{+} + \Phi^{\ast} + \cdots,
\qquad
\mathbf{v}
= \vec v + v^{+} + v^{\ast} + \cdots,
\qquad
\mathbf{S}
= S + S^{+} + S^{\ast} + \cdots,
\]
with ghost number assignments determined by the degree shifts of
$\mathsf{T}[1]M$ and $\mathrm{Plenum}$. The precise assignments follow
the standard AKSZ convention:
\[
\mathrm{gh}(\mathrm{field})=0,\qquad
\mathrm{gh}(\mathrm{ghost})=+1,\qquad
\mathrm{gh}(\mathrm{antifield})=-1,\qquad
\mathrm{gh}(\mathrm{higher~ghosts})>+1.
\]

\begin{proposition}
\label{prop:BV-square-zero}
The BV operator $Q_{\mathrm{BV}}$ is a degree $+1$ derivation satisfying
$Q_{\mathrm{BV}}^2=0$.
\end{proposition}

\begin{proof}[Proof (Sketch)]
$Q_{\mathrm{BV}}^2(F)=\{S,\{S,F\}\} = \frac12\{\{S,S\},F\}$, which
vanishes by Theorem~\ref{thm:CME-RSVP}.
\end{proof}


\section{Euler–Lagrange Equations and Lamphrodynamics}

We now extract the classical physical equations of motion as the
Euler–Lagrange equations of the AKSZ action at ghost number zero.

\begin{theorem}[RSVP Equations from AKSZ]
Let $S_{\mathrm{RSVP}}$ be the AKSZ action
\eqref{eq:RSVP-AKSZ}. Then the ghost-number-zero Euler--Lagrange
equations are equivalent to the lamphrodynamic field equations of
Chapter~5: namely,
\[
\Box \Phi = \lambda\nabla\!\cdot\vec{v} + \dots,
\qquad
\nabla\cdot \vec v = f(\nabla\Phi,\nabla S),
\qquad
\partial_t S + \nabla\!\cdot(S\vec v) = \kappa\nabla^2 S + \dots,
\]
with all nonlinear source terms determined by the same derived $3$–form
$\Theta$ used in the AKSZ construction.
\end{theorem}

\begin{proof}[Proof (Sketch)]
The variation of $S_{\mathrm{RSVP}}$ with respect to
$(\Phi,\vec v,S)$ at ghost number zero yields the classical field
equations. All antifield and ghost terms vanish at ghost number zero,
leaving precisely the lamphrodynamic PDE system of Chapter~5.
\end{proof}

\begin{remark}
Thus the AKSZ construction is not merely compatible with
lamphrodynamics; it {\em determines} the lamphrodynamic equations as the
classical sector of a $(-1)$–shifted symplectic BV theory. The derived
geometry forces the PDE system rather than assuming it.
\end{remark}


\section{Physical Interpretation}

The significance of the AKSZ construction is that it
endows RSVP with a canonical BV quantization framework without requiring
gauge-fixing or ad-hoc ghost terms: the ghost sector is automatic.
Moreover, the classical lamphrodynamic dynamics are precisely the
Euler--Lagrange equations of the BV action, making RSVP a field theory
in the strongest modern sense. Finally, the $(-1)$–shift ensures that
the BV bracket is well-defined and that the master equation holds,
a prerequisite for any rigorous quantum theory built on the BV
formalism.

\section{Boundary Data and BV--BFV Structures}
\label{sec:RSVP-BFV}

On a manifold with boundary $\partial M\neq\emptyset$, the AKSZ
construction produces a canonical BV--BFV structure in the sense of
Cattaneo--Mnev--Reshetikhin. Concretely, the boundary data
consist of a BFV space of boundary fields $\mathcal{F}_{\partial}$ and
a degree $0$ \emph{BFV action} $S_{\partial}$ which satisfies
\[
Q_{\mathrm{BFV}}^2=0,\qquad
Q_{\mathrm{BFV}}(S_{\partial})=0,
\]
and whose cohomology governs physical boundary degrees of freedom.

\begin{theorem}[Canonical BV--BFV Boundary Structure]
\label{thm:RSVP-BFV}
Let $M$ be a $4$--manifold with (possibly empty) boundary.
Then the AKSZ action $S_{\mathrm{RSVP}}$
induces a canonical BV--BFV structure
$(\mathcal{F}_{\partial},S_{\partial})$ on $\partial M$.
Moreover, $S_{\partial}$ determines the physical boundary excitations of
the lamphrodynamic fields $(\Phi,\vec v,S)$, including entropic flux
across $\partial M$.
\end{theorem}

\begin{proof}[Proof (Sketch)]
Immediate from the general BV--BFV theorem for AKSZ models, using that
the target is $(-1)$--shifted symplectic and that $M$ has dimension $4$.
The boundary inherits a degree-$0$ Hamiltonian system governing the
physical boundary data.
\end{proof}

\begin{remark}[Physical Meaning]
The presence of a canonical BV--BFV boundary theory implies that
entropic flux and vector circulation admit well-defined ``boundary''
degrees of freedom even before gauge-fixing.
This suggests a direct route to black-hole and horizon physics once
appropriate boundary conditions are chosen (Volume~II).
\end{remark}



\section{Local Coordinates and Expansion}
\label{sec:local-coords}

Let $(x^{\mu})_{\mu=0}^{3}$ be local coordinates on $M$ and
let $(\Phi,\vec v,S)$ denote the degree-$0$ components of the superfields.
Expanding the AKSZ action to leading (classical) order gives
\begin{equation}
\label{eq:local-RSVP}
S_{\mathrm{RSVP}} \simeq
\int_{M} \Big(
\partial_{\mu}\Phi \, \Theta^{\mu\nu\rho}
\, \partial_{\nu} v_{\rho}
+ \partial_{\mu}S \, \Theta^{\mu\nu\rho}
\, \partial_{\nu} \Phi
+ \dots
\Big)\, d^{4}x,
\end{equation}
where $\Theta^{\mu\nu\rho}$ are the components of the closed
$3$--form representing the $(-1)$--shifted symplectic structure
(cf. Theorem~\ref{thm:Main-Plenum}).

\begin{remark}
Subleading terms encode the nonlinear couplings responsible for
entropy flow and lamphrodynamic vorticity. These terms become crucial
in the strong-flux regime and near entropic singularities (Ch.~12).
\end{remark}



\section{Comparison to Chern--Simons AKSZ}
\label{sec:CS-comparison}

A useful analogy is provided by $3$--dimensional Chern--Simons theory,
which is the prototypical AKSZ model in dimension $3$ with target
$\mathfrak{g}[1]$.

\subsection{Chern--Simons AKSZ recap}
Let $N$ be a $3$--manifold and $\mathfrak{g}$ a Lie algebra.
The AKSZ construction gives the classical CS action
\[
S_{\mathrm{CS}} = \int_{N}
\Big( A\wedge dA + \tfrac{1}{3}A\wedge [A,A] \Big),
\]
arising from the canonical symplectic structure on $\mathfrak{g}[1]$
and the Chevalley--Eilenberg differential.

\subsection{Structural analogy}
Two structural parallels are immediate:

\begin{enumerate}
\item
{\bf Degree shift.}
CS uses a $0$--shifted target $\mathfrak{g}[1]$, whereas RSVP uses a
$(-1)$--shifted derived stack.
The lower shift reflects the presence of physical scalar
and vector fields rather than a pure gauge sector.

\item
{\bf Closed form generating the action.}
In both theories, the action is generated from the geometric
data of the target:
the Lie algebra $(\mathfrak{g},[-,-])$ in CS,
the lamphrodynamic $3$--form $\Theta$ in RSVP.
\end{enumerate}

\subsection{Physical meaning}
For Chern--Simons, the AKSZ construction explains why the CS action is
automatically topological and why its BV–BFV boundary theory captures
edge modes (e.g.\ quantum Hall).
In RSVP the AKSZ construction explains why lamphrodynamic equations
arise as Euler--Lagrange equations of a $(-1)$--shifted BV theory,
and thus why entropy flow and vector circulation are
\emph{geometric} rather than phenomenological.

\begin{remark}
The fact that lamphrodynamics sits one degree below the CS example
suggests a natural comparison between topological edge modes and
entropy-flow boundary modes. This will reappear in Chapter~12
(entropic singularities) and Volume~II (black-hole analogs).
\end{remark}



\section{Summary and Forward Pointer}
\label{sec:summary-aksz}

In this chapter we have:
\begin{enumerate}
\item Constructed the AKSZ action for RSVP from the
      $(-1)$--shifted target $\mathrm{Plenum}$.
\item Verified the classical master equation.
\item Identified the BV operator and ghost/antifield expansion.
\item Derived lamphrodynamic PDEs as Euler--Lagrange equations.
\item Described canonical boundary BFV structure.
\item Compared RSVP to the Chern--Simons AKSZ example.
\end{enumerate}

In the next chapter we develop the variational and Hamiltonian
formulation of RSVP in classical degree zero, proving that the
$(-1)$--shifted symplectic geometry induces a well-defined Hamiltonian
system and enabling energy estimates and Noether-type conservation
laws. This prepares the ground for the analytic theory developed in
Chapters~9--11.


\chapter{Variational Principles and Hamiltonian Formulation}
\label{ch:variational}

The AKSZ construction of the previous chapter provides a BV
master action $S_{\mathrm{RSVP}}$ on the graded space of superfields.
In this chapter we extract the classical (degree~$0$) variational
principle, identify the induced Euler--Lagrange equations in ordinary
fields $(\Phi,\vec v,S)$, and construct the classical Hamiltonian
formulation on the derived mapping stack.
This leads to a well-defined notion of conserved currents and prepares
the ground for the analytic theory developed in Chapters~9--11.



\section{The Classical Variational Functional}
\label{sec:variational-functional}

Let $M$ be a smooth, oriented $4$--manifold without boundary (the
boundary case will be handled in the BV--BFV formalism).
Restrict the AKSZ action $S_{\mathrm{RSVP}}$ to the classical (ghost
number~$0$) sector. Denote by
\[
(\Phi,\vec v,S)\in \Gamma(M,\mathcal{E})
\]
the physical lamphrodynamic fields.
Then the classical action takes the form
\begin{equation}
\label{eq:classical-action}
\mathcal{S}[\Phi,\vec v,S]
=\int_{M}\mathcal{L}(\Phi,\vec v,S,\partial\Phi,\partial\vec v,
\partial S)\,d^{4}x,
\end{equation}
where the Lagrangian density $\mathcal{L}$ is obtained by truncating
the AKSZ superfield expansion and selecting the degree~$0$ component
(including nonlinear terms arising from the $(-1)$--shifted
symplectic structure).

\begin{remark}
The explicit coordinate expression of $\mathcal{L}$ is cumbersome due
to the nonlinear couplings between $\Phi,\vec v$ and $S$. However, the
structural properties---gauge invariances, conservation laws, and
entropy-flux terms---follow from the $(-1)$--shifted geometry rather
than from the coordinate expansion.
\end{remark}



\section{Euler--Lagrange Equations}
\label{sec:EL-equations}

The Euler--Lagrange equations are obtained by taking first variations
of \eqref{eq:classical-action}:
\begin{align}
\frac{\delta\mathcal{S}}{\delta\Phi}
&=
\partial_{\mu}\left(
\frac{\partial\mathcal{L}}{\partial(\partial_{\mu}\Phi)}
\right)
-
\frac{\partial\mathcal{L}}{\partial\Phi}
=0,
\label{eq:EL-Phi}
\\[0.5em]
\frac{\delta\mathcal{S}}{\delta v_{\nu}}
&=
\partial_{\mu}\left(
\frac{\partial\mathcal{L}}{\partial(\partial_{\mu}v_{\nu})}
\right)
-
\frac{\partial\mathcal{L}}{\partial v_{\nu}}
=0,
\label{eq:EL-v}
\\[0.5em]
\frac{\delta\mathcal{S}}{\delta S}
&=
\partial_{\mu}\left(
\frac{\partial\mathcal{L}}{\partial(\partial_{\mu}S)}
\right)
-
\frac{\partial\mathcal{L}}{\partial S}
=0.
\label{eq:EL-S}
\end{align}

\begin{theorem}[Equivalence with AKSZ Euler--Lagrange Equations]
\label{thm:EL-equivalence}
The equations \eqref{eq:EL-Phi}--\eqref{eq:EL-S} coincide with the
classical Euler--Lagrange equations obtained from the AKSZ action
$S_{\mathrm{RSVP}}$ by projecting to degree~$0$.
\end{theorem}

\begin{proof}
The AKSZ construction guarantees that the BV Euler--Lagrange
operator $\delta S_{\mathrm{RSVP}}/\delta\mathsf{Field}=0$ restricts
to the classical sector by ghost number.
Truncating to ghost number~$0$ and using the fact that the BV bracket
reduces to the classical Poisson bracket on physical fields yields
\eqref{eq:EL-Phi}--\eqref{eq:EL-S}.
\end{proof}



\section{Hamiltonian Formulation on the Derived Mapping Stack}
\label{sec:Hamiltonian}

Let $\mathrm{Map}(M,\mathrm{Plenum})$ be the classical mapping space
of $M$ into the derived target $\mathrm{Plenum}$,
and let $\omega_{(-1)}$ denote the $(-1)$--shifted symplectic form on
$\mathrm{Plenum}$.
In the AKSZ formalism, the Hamiltonian structure on the space of
superfields descends (after truncation) to a classical Hamiltonian
system on $\mathrm{Map}(M,\mathrm{Plenum})$ endowed with the induced
(shifted) 2--form.

\begin{proposition}[Classical Hamiltonian Structure]
\label{prop:Hamiltonian}
The classical space of fields $\mathrm{Map}(M,\mathrm{Plenum})$
carries a well-defined Hamiltonian structure $(\omega,H)$ where
\[
H=\int_{M} \mathcal{H}(\Phi,\vec v,S,
\partial\Phi,\partial\vec v,\partial S)\,d^{4}x,
\]
and the Hamiltonian density~$\mathcal{H}$ is obtained by classical
Legendre transform of the Lagrangian density~$\mathcal{L}$.
\end{proposition}

\begin{proof}[Proof (Sketch)]
The $(-1)$--shifted symplectic form on $\mathrm{Plenum}$ induces a
$0$--shifted (classical) symplectic form on the mapping space.
The Legendre transform is defined fiberwise on the classical field
bundle $\mathcal{E}\to M$.
Standard AKSZ arguments show that the resulting Hamiltonian generates
the same equations of motion as \eqref{eq:EL-Phi}--\eqref{eq:EL-S}.
\end{proof}

\begin{remark}[Odd vs Even Poisson Brackets]
The BV bracket on the AKSZ superfields is odd (degree $+1$),
whereas the induced bracket on the classical fields is even.
This provides a natural separation between quantum BV information
(ghost/antifield sector) and classical conservation laws.
\end{remark}

\section{Noether Currents and Conserved Quantities}
\label{sec:noether}

Let $\mathcal{G}$ be a Lie group of symmetries acting on the classical
fields $(\Phi,\vec v,S)$ via smooth bundle automorphisms
compatible with the lamphrodynamic constraint.
Write $\rho:\mathfrak{g}\to\Gamma(T\mathcal{E})$ for the induced Lie
algebra action on fields.

\begin{definition}[Infinitesimal Symmetry]
A vector field $X\in\Gamma(T\mathcal{E})$ is an infinitesimal
symmetry if the Lie derivative of the Lagrangian density satisfies
\[
\mathcal{L}_{X}\mathcal{L} = \partial_{\mu} B^{\mu}
\]
for some current $B^{\mu}$.
\end{definition}

Standard variational arguments yield:

\begin{proposition}[Noether Current]
\label{prop:noether-current}
Let $X$ be an infinitesimal symmetry associated with $\xi\in\mathfrak{g}$.
Then the Noether current is
\[
J^{\mu}_{\xi}
= \frac{\partial\mathcal{L}}{\partial(\partial_{\mu}\Phi)}X(\Phi)
 +\frac{\partial\mathcal{L}}{\partial(\partial_{\mu}v_{\nu})}X(v_{\nu})
 +\frac{\partial\mathcal{L}}{\partial(\partial_{\mu}S)}X(S)
 - B^{\mu},
\]
and is conserved along classical solutions:
$\partial_{\mu}J^{\mu}_{\xi}=0$.
\end{proposition}

\begin{remark}
In the BV formulation of Chapter~7, the same current arises as the
degree--$0$ component of the BV charge associated to the ghost--$\xi$.
Hence classical conservation laws are shadows of BV identities.
\end{remark}



\section{Hamiltonian Vector Fields and Poisson Structure}
\label{sec:ham-vfields}

Let $\omega$ denote the classical symplectic form on
$\mathrm{Map}(M,\mathrm{Plenum})$.
For any classical observable $\mathcal{O}$, define the Hamiltonian
vector field $X_{\mathcal{O}}$ by the relation
\[
\iota_{X_{\mathcal{O}}}\omega = \delta\mathcal{O}.
\]
Provided $\omega$ is weakly nondegenerate on the classical sector,
this determines $X_{\mathcal{O}}$ formally and induces a Poisson
bracket
\[
\{\mathcal{O}_{1},\mathcal{O}_{2}\}
:= X_{\mathcal{O}_{1}}(\mathcal{O}_{2}),
\]
which restricts to field polynomials under classical truncation.

\begin{remark}
Weak nondegeneracy follows from nondegeneracy of the $(-1)$--shifted
symplectic form on the target stack $\mathrm{Plenum}$ and the AKSZ
descent; a detailed discussion will appear in Appendix~E.
\end{remark}



\section{Energy, Momentum, and Entropic Contributions}
\label{sec:energy-momentum}

Energy--momentum for lamphrodynamic fields requires special care
because $S$ contributes an entropic stress--energy component not
present in standard relativistic field theories.

\begin{definition}[Lamphrodynamic Stress--Energy Tensor]
\label{def:stress}
Define the symmetric tensor
\[
T_{\mu\nu}
:=
\frac{\partial\mathcal{L}}{\partial(\partial^{\mu}\Phi)}
 \partial_{\nu}\Phi
+\frac{\partial\mathcal{L}}{\partial(\partial^{\mu}v_{\rho})}
 \partial_{\nu} v_{\rho}
+\frac{\partial\mathcal{L}}{\partial(\partial^{\mu}S)}
 \partial_{\nu}S
- g_{\mu\nu}\mathcal{L},
\]
where $g_{\mu\nu}$ is the background Lorentzian metric on $M$.
\end{definition}

\begin{proposition}[Energy Balance Identity]
\label{prop:energy-balance}
Along classical solutions of the Euler--Lagrange equations
\eqref{eq:EL-Phi}--\eqref{eq:EL-S},
the stress--energy tensor satisfies
\[
\nabla^{\mu}T_{\mu\nu}
= \mathscr{E}_{\nu}(\Phi,\vec v,S),
\]
where the entropic source term
$\mathscr{E}_{\nu}$ vanishes for isentropic flows and encodes
entropy--flux contributions in general.
\end{proposition}

\begin{remark}
The presence of $\mathscr{E}_{\nu}$ is physically crucial: it
generates effective geometric backreaction in emergent gravity
(Volume~II, Chapter~1).
\end{remark}



\section{Hamiltonian Energy Estimates and PDE Well--Posedness}
\label{sec:ham-energy}

The classical Hamiltonian
$H=\int_{M}\mathcal{H}\,d^{4}x$
provides the fundamental energy control for the analytic theory of
Chapters~9--11.
The following preparatory result connects the formal Hamiltonian
with Sobolev control.

\begin{theorem}[Hamiltonian Energy Estimate]
\label{thm:ham-estimate}
Let $(\Phi,\vec v,S)$ be a classical solution of the lamphrodynamic
field equations on a globally hyperbolic spacetime $(M,g)$.
Assume initial data in $H^{k}$, $k\ge2$,
and suppose the entropic term satisfies a coercivity condition
$\mathscr{E}_{0}\le C|\nabla S|^{2}$.
Then for sufficiently small time $T>0$,
\[
\sup_{t\in[0,T]}
\bigl( \|\Phi(t)\|_{H^{k}}
     +\|\vec v(t)\|_{H^{k}}
     +\|S(t)\|_{H^{k}}
\bigr)
\;\le\;
C_{0} \exp(CT)\,
\bigl( \|\Phi_{0}\|_{H^{k}}
     +\|\vec v_{0}\|_{H^{k}}
     +\|S_{0}\|_{H^{k}} \bigr),
\]
where $C_{0},C$ depend on the geometry of $M$.
\end{theorem}

\begin{proof}[Proof (Sketch)]
The Hamiltonian controls the $H^{k}$--norms modulo entropic terms.
Coercivity of $\mathscr{E}_{0}$ and Gr\"onwall's inequality yield the
exponential bound.
A detailed PDE proof is given in Chapter~10.
\end{proof}



\section{Summary and Forward Link}
\label{sec:summary}

In summary, the classical limit of the AKSZ--RSVP theory yields:
\begin{enumerate}
\item a well-defined variational principle,
\item Euler--Lagrange equations equivalent to the classical AKSZ equations,
\item a Hamiltonian formalism on the derived mapping stack,
\item Noether currents and entropic stress--energy,
\item energy estimates sufficient for local well-posedness.
\end{enumerate}
These results serve as the bridge between the geometric BV framework
(Chapters~4--7) and the analytic theory (Chapters~9--11),
where energy estimates and Sobolev bounds play a central role.



% ------------------------------------------------------------
% Part III: Well-Posedness and Analysis
% ------------------------------------------------------------
\part{Well-Posedness and Analysis}

\chapter{Linear Theory}
\label{ch:linear}

Chapters~5--8 developed the geometric and variational structure of
lamphrodynamic fields $(\Phi,\vec v,S)$, culminating in the classical
Hamiltonian formulation.
We now turn to analytic questions, beginning with the linearization of
the Euler--Lagrange equations \eqref{eq:EL-Phi}--\eqref{eq:EL-S} about
a smooth reference configuration and establishing global well--posedness
in Sobolev spaces.
This chapter collects the functional--analytic tools needed for the
nonlinear theory of Chapters~10--11.



\section{Linearization about a Reference Configuration}
\label{sec:linearization}

Let $(\Phi,\vec v,S)$ be a smooth lamphrodynamic configuration solving
the classical Euler--Lagrange equations on a globally hyperbolic
spacetime $(M,g)$.
Define the linearized fields by
\[
(\phi, w_{\mu}, s)
:=
\frac{d}{d\varepsilon}\big|_{\varepsilon=0}
(\Phi+\varepsilon\phi,\; v_{\mu}+\varepsilon w_{\mu},\; S+\varepsilon s).
\]
To simplify notation, write $U:=(\phi, w, s)$.
Substituting into \eqref{eq:EL-Phi}--\eqref{eq:EL-S} and retaining
terms linear in $U$ yields a linear system of wave--type operators
\begin{equation}
\label{eq:linear-system}
\mathcal{L}_{\mathrm{lin}} U
= \square_{g} U
 + B^{\mu}\nabla_{\mu}U
 + C\,U
=0,
\end{equation}
where $\square_{g}$ is the tensor wave operator, $B^{\mu}$ is a
first--order coefficient determined by the background
$(\Phi,\vec v,S)$, and $C$ is a zero--order coefficient.
The coefficients $B^{\mu},C$ are smooth and bounded on compact sets.



\begin{remark}[Hyperbolic Character]
The principal part of \eqref{eq:linear-system} is the wave operator,
ensuring hyperbolicity; lower--order terms encode entropy--flux
corrections of the background.
\end{remark}



\section{Functional Spaces and Initial Value Problem}
\label{sec:spaces}

Fix a spacelike Cauchy hypersurface~$\Sigma$ with induced Riemannian
metric~$h$.
For integer $k\ge0$ define Sobolev norms
\[
\|U\|_{H^{k}(\Sigma)}
 :=
 \sum_{j=0}^{k}
  \| \nabla^{j} U \|_{L^{2}(\Sigma)}.
\]

\begin{definition}[Initial Value Problem]
\label{def:IVP}
Given initial data
\[
U(0)=U_{0}\in H^{k}(\Sigma),\qquad
\partial_{t}U(0)=U_{1}\in H^{k-1}(\Sigma),
\]
find $U(t,x)$ solving \eqref{eq:linear-system} for $t$ in some time
interval containing $0$.
\end{definition}



\section{Energy Estimates}
\label{sec:linear-energy}

Define the linear energy functional
\begin{equation}
\label{eq:linear-energy}
E_{\mathrm{lin}}(U)(t)
:= \frac{1}{2}
\int_{\Sigma_{t}}
\bigl(
 |\partial_{t}U|^{2}
 +|\nabla U|^{2}
 + |U|^{2}
\bigr)\, d\mathrm{vol}_{h}.
\end{equation}
Standard integration--by--parts yields:

\begin{lemma}[Linear Energy Estimate]
\label{lem:linear-energy}
For sufficiently smooth $U$ solving \eqref{eq:linear-system},
\[
\frac{d}{dt}E_{\mathrm{lin}}(U)(t)
\le
C\bigl(E_{\mathrm{lin}}(U)(t)+\|U(t)\|_{L^{2}}^{2}\bigr),
\]
where $C$ depends on the background coefficients.
\end{lemma}

\begin{proof}[Proof (Sketch)]
Multiply \eqref{eq:linear-system} by $\partial_{t}U$,
integrate over $\Sigma_{t}$,
use divergence theorem on the wave operator, and absorb
lower--order terms into the right--hand side.
\end{proof}

By Grönwall's inequality we obtain:

\begin{corollary}[A priori Control]
\label{cor:apriori}
For $t\in[0,T]$,
\[
E_{\mathrm{lin}}(U)(t)
\le
E_{\mathrm{lin}}(U)(0)\,\exp(C T).
\]
\end{corollary}



\section{Global Well--posedness}
\label{sec:global-wp}

We now state the main linear result.

\begin{theorem}[Global Well--posedness in Sobolev Spaces]
\label{thm:linear-wp}
Let $(M,g)$ be globally hyperbolic and let $k\ge2$.
For any initial data $(U_{0},U_{1})\in H^{k}(\Sigma)\times H^{k-1}(\Sigma)$
there exists a unique global solution
\[
U\in C^{0}([0,\infty);H^{k}(\Sigma))
 \cap C^{1}([0,\infty);H^{k-1}(\Sigma))
\]
of \eqref{eq:linear-system}.
Moreover, the energy bound of Corollary~\ref{cor:apriori} holds.
\end{theorem}

\begin{proof}[Proof]
Local existence and uniqueness follow from standard hyperbolic PDE
theory for systems with smooth coefficients
\cite{Hormander-AnalysisPDE}.
Global extension follows from the a priori bound
(Corollary~\ref{cor:apriori}) combined with continuation criteria.
\end{proof}



\section{Spectral Properties of the Linearized Operator}
\label{sec:spectral}

On static spacetimes or after Fourier transform in time,
$\mathcal{L}_{\mathrm{lin}}$ reduces to an elliptic operator of the
form
\[
\mathscr{A} = -\Delta_{h} + V(x),
\]
where $V$ is a potential induced by the background.
Standard spectral theory implies:

\begin{proposition}[Spectral Properties]
\label{prop:spectral}
If $(M,g)$ is static and $\Sigma$ compact, then
$\mathscr{A}$ has discrete spectrum $\{\lambda_{j}\}_{j\ge0}$,
with $\lambda_{j}\to+\infty$ as $j\to\infty$.
\end{proposition}

\begin{proof}
Immediate from ellipticity and compactness of $\Sigma$
\cite{Taylor-PDE}.
\end{proof}



\section{Connection to AKSZ Classical Limit}
\label{sec:connection-AKSZ-classical}

The linear theory developed here governs fluctuations around classical
solutions of the AKSZ--RSVP action. The coercivity of the Hamiltonian
(Chapter~8) and the linear estimate above ensure that the BV
perturbative expansion is formally well behaved around smooth
backgrounds, a prerequisite for quantization in Volume~III.



\section{Summary}

We have shown that linearized lamphrodynamic fields satisfy a
hyperbolic system of second--order PDEs with global well--posedness
in Sobolev spaces. The resulting energy and spectral control forms the
analytic basis for the nonlinear theory of Chapters~10--11.


\chapter{Energy Estimates and Regularity}
\label{ch:energy-regularity}

The global well--posedness of the linearized lamphrodynamic system
(Chapter~\ref{ch:linear}) forms the analytic foundation for the
nonlinear theory.  We now establish higher--order energy estimates for
the full nonlinear lamphrodynamic equations, prove local well--posedness
in Sobolev spaces, and derive regularity criteria that exert control
near potential entropic singularities.  These results prepare the ground
for the global continuation problem in Chapter~11.



\section{Nonlinear Lamphrodynamic System}
\label{sec:nonlinear-system}

Let $(\Phi,\vec v,S)$ satisfy the Euler--Lagrange equations
\eqref{eq:EL-Phi}--\eqref{eq:EL-S}.
Suppressing indices, write the nonlinear system schematically as
\begin{equation}
\label{eq:nonlinear-system}
\square_{g}U
=
\mathcal{N}(U,\nabla U),
\end{equation}
where $U=(\Phi,\vec v,S)$ and $\mathcal{N}$ is a smooth nonlinear
mapping whose algebraic form depends polynomially on $U$ and its first
derivatives; in particular, no derivative loss occurs at the level of
principal part.



\section{Higher--order Energies}
\label{sec:higher-order-energies}

For integer $k\ge2$ define the $k$--th order energy
\begin{equation}
\label{eq:Ek}
E_{k}(t)
=
\frac{1}{2}\!
\sum_{|\alpha|\le k}
\int_{\Sigma_{t}}
\Bigl(
|\partial_{t}\nabla^{\alpha}U|^{2}
 + |\nabla\nabla^{\alpha}U|^{2}
 + |\nabla^{\alpha}U|^{2}
\Bigr)\, d\mathrm{vol}_{h},
\end{equation}
where $\alpha$ is a spatial multi-index and $\nabla^{\alpha}$ denotes
the covariant derivative with respect to $h$.

Our goal is to estimate $E_{k}(t)$ in terms of initial data
$E_{k}(0)$ and lower--order terms.



\begin{lemma}[Commutator Estimate]
\label{lem:commutator}
For each $|\alpha|\le k$ one has
\[
\bigl\|\![\nabla^{\alpha},\mathcal{N}]\,U\bigr\|_{L^{2}(\Sigma_{t})}
\le
C\sum_{j\le k}\|\nabla^{j}U\|_{L^{2}(\Sigma_{t})}^{2}.
\]
\end{lemma}

\begin{proof}[Proof (Sketch)]
The nonlinearities are polynomial in $U$ and $\nabla U$ with smooth
coefficients, so commutators produce lower--order terms which are
controlled by standard Moser--type estimates.
\end{proof}



\section{Nonlinear Energy Inequality}
\label{sec:nonlinear-energy}

Commuting $\nabla^{\alpha}$ with \eqref{eq:nonlinear-system}, multiplying
by $\partial_{t}\nabla^{\alpha}U$, integrating over $\Sigma_{t}$, and
using Lemma~\ref{lem:commutator} yields the fundamental inequality
\begin{equation}
\label{eq:nonlinear-energy-ineq}
\frac{d}{dt}E_{k}(t)
\le
C\bigl(E_{k}(t) + E_{k}(t)^{3/2}\bigr).
\end{equation}

\begin{remark}
The cubic growth reflects the polynomial nonlinearity of the
lamphrodynamic interactions; the crucial point is absence of derivative
loss, ensuring that the $H^{k}$--norm remains the natural energy scale.
\end{remark}



\section{Local Well--posedness}
\label{sec:local-wp}

We now state the main nonlinear local existence theorem.

\begin{theorem}[Local Well--posedness in $H^{k}$]
\label{thm:local-wp}
Let $k\ge2$ and initial data
$U_{0}\in H^{k}(\Sigma)$,
$U_{1}\in H^{k-1}(\Sigma)$.
Then there exists $T>0$ and a unique solution
\[
U\in C^{0}\bigl([0,T];H^{k}(\Sigma)\bigr)
\cap C^{1}\bigl([0,T];H^{k-1}(\Sigma)\bigr)
\]
of the nonlinear system \eqref{eq:nonlinear-system}.
\end{theorem}

\begin{proof}[Proof (Sketch)]
Iterate the linear problem with nonlinear forcing; apply the
energy inequality \eqref{eq:nonlinear-energy-ineq} and Gr\"onwall’s
inequality to obtain uniform bounds on a short time interval.  Contraction
mapping in the energy norm yields uniqueness.
\end{proof}



\section{Continuation Criterion and Blowup Scenarios}
\label{sec:continuation}

Let $T_{\max}$ denote the maximal time of existence in the class of
$H^{k}$--solutions.

\begin{theorem}[Continuation Criterion]
\label{thm:continuation}
If $T_{\max}<\infty$, then
\[
\limsup_{t\nearrow T_{\max}}
\Bigl(
\|U(t)\|_{H^{k}}
+\|\nabla U(t)\|_{L^{\infty}}
\Bigr)
=+\infty.
\]
\end{theorem}

\begin{proof}[Proof (Sketch)]
Suppose the right--hand side remains finite; controlling $\nabla U$ in
$L^{\infty}$ yields uniform bounds for nonlinear terms, contradicting
maximality.  Standard hyperbolic PDE techniques apply.
\end{proof}

\begin{remark}[Possible Singularities]
Entropic singularities may arise precisely when gradients of $S$
become unbounded.  Understanding such scenarios is central to the global
existence problem (Chapter~11).
\end{remark}



\section{Higher Regularity and Sobolev Embeddings}
\label{sec:higher-regularity}

Assume the solution $U$ lies in $H^{k}$ for $k>\frac{n}{2}+1=3$ in
spatial dimension~$3$; then Sobolev embedding implies
\[
\|U\|_{C^{1}} \le C\|U\|_{H^{k}}.
\]
Together with \eqref{eq:nonlinear-energy-ineq}, this yields:
\begin{corollary}[Higher Regularity]
\label{cor:higher-reg}
If initial data belong to $H^{k}$, $k>3$, then the solution remains
$C^{1}$ for as long as it exists. \qed
\end{corollary}



\section{Entropy Coercivity and Stability}
\label{sec:entropy-coercive}

Entropy plays a stabilizing role when its contribution satisfies a
coercivity condition similar to that in
Theorem~\ref{thm:ham-estimate}.
Assume
\begin{equation}
\label{eq:entropy-coercive}
\mathscr{E}_{0}(U)\le C|\nabla S|^{2}.
\end{equation}
Then:

\begin{proposition}[Stability under Entropy Coercivity]
\label{prop:stability}
If \eqref{eq:entropy-coercive} holds, then there exist constants
$C_{0},C>0$ such that
\[
E_{k}(t) \le C_{0}E_{k}(0)\exp(CT),
\qquad 0\le t\le T.
\]
\end{proposition}

\begin{proof}[Proof (Sketch)]
Combine \eqref{eq:entropy-coercive} with
\eqref{eq:nonlinear-energy-ineq}
and apply Gr\"onwall’s inequality.
\end{proof}



\section{Summary and Forward Link}
\label{sec:summary10}

We have proved local well--posedness, higher--order energy estimates,
and a continuation criterion for the full nonlinear lamphrodynamic
equations.  The analysis hinges on coercivity of entropic terms and
the absence of derivative loss.  In Chapter~\ref{ch:global}, we study
global continuation, entropy--driven singularities, and prospects for
global existence.


\chapter{Nonlinear Global Theory}
\label{ch:global}

The local theory of Chapter~\ref{ch:energy-regularity} guarantees
existence and uniqueness of classical lamphrodynamic solutions on
a short time interval, subject to finite Sobolev norms.
Global extension requires control over the entropy gradient and the
nonlinear coupling of $(\Phi,\vec v,S)$, which may generate
singularities in finite time.
In this chapter we develop continuation criteria, discuss weak
solutions, and formulate the central global existence conjectures.
Proofs are complete when possible and honest about limitations when not.



\section{Entropy--Driven Singularities}
\label{sec:entropy-sing}

Let $U=(\Phi,\vec v,S)$ solve the nonlinear system
\eqref{eq:nonlinear-system}.
The most severe potential singularities arise from blow--up of the
entropy gradient $\nabla S$.

\begin{proposition}[Entropy Singularity Criterion]
\label{prop:entropy-sing}
If a classical solution loses smoothness at finite time
$T_{\max}<\infty$, then
\[
\limsup_{t\nearrow T_{\max}}
\|\nabla S(t)\|_{L^{\infty}} = +\infty.
\]
\end{proposition}

\begin{proof}[Proof (Sketch)]
If $\|\nabla S(t)\|_{L^{\infty}}$ remained bounded, then by
Theorem~\ref{thm:continuation}, the $H^{k}$--norms would also remain
bounded, contradicting maximality of $T_{\max}$.
\end{proof}

\begin{remark}
This establishes the entropy gradient as the natural blow--up variable
for lamphrodynamic fields, in contrast to standard fluid models where
density or vorticity dominate.
\end{remark}



\section{Global Existence under Small Data}
\label{sec:small-data}

On Minkowski space $\mathbb{R}^{1+3}$, global existence can be shown
for sufficiently small initial data by exploiting the dispersive
character of the wave operator.

\begin{theorem}[Small--Data Global Existence]
\label{thm:small-data}
Let $M=\mathbb{R}^{1+3}$ with flat metric and assume
$\|U_{0}\|_{H^{k}}+\|U_{1}\|_{H^{k-1}}$ is sufficiently small.
Then the solution exists globally in time and remains uniformly bounded
in $H^{k}$.
\end{theorem}

\begin{proof}[Proof (Sketch)]
Use dispersive decay for the linear wave equation and perturbative
arguments controlling the nonlinear term by the dispersive gain.
The absence of derivative loss is essential here.
\end{proof}



\section{Weak Solutions and Entropy Inequalities}
\label{sec:weak-sol}

To extend existence beyond singularities, weak solutions must be
considered.  Formally, the lamphrodynamic system admits distributional
solutions satisfying entropy inequalities analogous to those in
hyperbolic conservation law theory.

\begin{definition}[Weak Solution]
A function $U\in L^{2}_{\mathrm{loc}}(M)$ is a weak solution if
\eqref{eq:nonlinear-system} holds in the sense of distributions and
$U$ satisfies an entropy inequality
\[
\partial_{t}\eta(U)
+\nabla\!\cdot\! q(U)
\le 0
\]
for all convex entropy--entropy flux pairs $(\eta,q)$.
\end{definition}

\begin{remark}
The entropy inequality imposes an additional physical admissibility
constraint consistent with the variational structure and prevents
unphysical shocks in $S$.
\end{remark}



\section{Global Existence in the Weak Class}
\label{sec:weak-global}

Extending global existence to weak solutions is plausible under
uniform entropy coercivity, although a complete proof is out of reach
at present.

\begin{conjecture}[Weak Global Existence]
\label{conj:weak-global}
Let $(M,g)$ be globally hyperbolic and assume the entropy coercivity
condition~\eqref{eq:entropy-coercive}.
Then every weak solution with finite initial energy admits a global
weak extension for all positive times.
\end{conjecture}

\begin{remark}
Conjecture~\ref{conj:weak-global} parallels global weak existence in
incompressible Euler but is complicated by entropy dynamics and
derived geometric structure.
\end{remark}



\section{Nonlinear Stability of Smooth Backgrounds}
\label{sec:stability}

Let $U_{\mathrm{bg}}$ be a smooth global solution (``background state'').
Define perturbations $u:=U-U_{\mathrm{bg}}$.

\begin{theorem}[Nonlinear Stability (conditional)]
\label{thm:nonlinear-stability}
Assume $U_{\mathrm{bg}}$ is globally smooth and satisfies uniform
entropy coercivity.  Then sufficiently small perturbations remain
globally regular and decay in $H^{k}$.
\end{theorem}

\begin{proof}[Proof (Sketch)]
Linear decay around the background combined with bootstrap bounds for
nonlinear terms yields uniform control provided the entropy coercivity
remains effective.
\end{proof}

\begin{remark}
A full proof requires refined analysis of the background state, omitted
here.  This remains an important research direction (see Volume~III).
\end{remark}



\section{Global Existence vs.~Entropic Singularity}
\label{sec:dichotomy}

Combining Proposition~\ref{prop:entropy-sing},
Theorem~\ref{thm:small-data},
and Conjecture~\ref{conj:weak-global}, we obtain the following
conceptual dichotomy:

\begin{theorem}[Global Dichotomy]
\label{thm:dichotomy}
Every maximal classical solution exhibits one of the following:
\begin{enumerate}
\item global classical existence, or
\item finite--time blow--up of the entropy gradient.
\end{enumerate}
\end{theorem}

\begin{proof}
If the solution does not remain classical globally, then
Proposition~\ref{prop:entropy-sing} forces blow--up of $\nabla S$.
\end{proof}



\section{Open Problems}
\label{sec:open-problems-11}

\begin{enumerate}
\item \textbf{Global classical existence.}
Does entropy coercivity suffice to prevent blow--up of $\nabla S$
for generic initial data?

\item \textbf{Weak--strong uniqueness.}
If a classical solution exists, must any weak solution coincide with it?

\item \textbf{Entropy singularity formation.}
Can entropic singularities form from smooth, finite--energy data?

\item \textbf{Critical regularity theory.}
Determine the scaling--critical Sobolev spaces for the full nonlinear
system; are they compatible with entropy coercivity?

\item \textbf{Numerical evidence.}
Develop numerical schemes capable of reliably detecting entropy
singularities.
\end{enumerate}



\section{Summary}

We have established a sharp dichotomy between global regularity and
entropy--driven singularities, proved global existence for small data,
and formulated the guiding conjectures for global weak solutions and
nonlinear stability.  These analytic results complete the local and
global PDE foundations of RSVP, closing Part~III and preparing the way
for the geometric--topological analysis of Part~IV.



% ------------------------------------------------------------
% Part IV: Geometric and Topological Properties
% ------------------------------------------------------------
\part{Geometric and Topological Properties}

\chapter{Entropic Singularities and Topology}
\label{ch:entropic-topology}

The analysis of Chapter~\ref{ch:global} shows that failure of classical
regularity is driven by blow--up of the entropy gradient $\nabla S$.
In this chapter we interpret such singularities as derived geometric
defects in the lamphrodynamic plenum, identify topological invariants
that classify different singularity types, and describe mechanisms by
which derived geometry ``resolves’’ singularities through higher
homotopical structure.
This constitutes the geometric completion of the analytic story.



\section{Derived Singularities of the Plenum Stack}
\label{sec:derived-sing}

Let
\[
\mathrm{Plenum}
= \mathcal{U}
\]
be the derived Artin stack of lamphrodynamic configurations constructed
in Chapter~\ref{ch:plenum-construction}.  By
Theorem~\ref{thm:Main-Plenum}, $\mathrm{Plenum}$ carries a
$2$--shifted symplectic structure whose AKSZ descent produces the
classical BV theory of Chapter~7.

\begin{definition}[Derived Entropic Singularity]
A point $x\in M$ is an \emph{entropic singularity} of a classical
configuration if
\[
|\nabla S(x)|=+\infty.
\]
In the derived picture, the corresponding point in
$\mathrm{Map}(M,\mathrm{Plenum})$ lies outside the classical truncation
but admits a well--defined point in the derived mapping stack.
\end{definition}

\begin{remark}
Derived geometry provides an intrinsic meaning to singular field
configurations: they exist as derived points even when no classical
representative exists.  This is why lamphrodynamics “continues’’
through singularities (cf.~Conjecture~\ref{conj:weak-global}).
\end{remark}



\section{Classification of Singularities}
\label{sec:classification}

Let $U=(\Phi,\vec v,S)$ be classical away from a closed set
$\Sigma_{\mathrm{sing}}\subset M$.
The induced map
\[
U\colon M\setminus\Sigma_{\mathrm{sing}}\longrightarrow \mathrm{Plenum}
\]
extends to a derived morphism
$\widetilde U\colon M\to\mathrm{Plenum}$.
Singularities thus correspond to nontrivial derived fibers at points of
$\Sigma_{\mathrm{sing}}$.

\begin{proposition}[Topological Classification]
\label{prop:classification}
Connected components of $\Sigma_{\mathrm{sing}}$ are classified up to
homotopy by elements of $\pi_{*}(\mathrm{Plenum})$, i.e.\ by homotopy
groups of the lamphrodynamic plenum.
\end{proposition}

\begin{proof}[Proof (Sketch)]
The derived extension $\widetilde U$ replaces the omission of
$\Sigma_{\mathrm{sing}}$ by higher homotopical data in
$\mathrm{Plenum}$, so homotopy classes of defects correspond to
homotopy classes of maps into the target.
\end{proof}

\begin{remark}
This phenomenon parallels topological defects in gauge theory (monopoles,
vortices), but here the relevant target is not a Lie group but the
derived stack $\mathrm{Plenum}$, producing a richer stratification of
defect types.
\end{remark}



\section{Characteristic Classes and Entropy}
\label{sec:char-classes}

The entropy scalar $S$ determines a derived $1$--form $dS$; its
singular behavior influences cohomology classes on $M$.

\begin{definition}[Entropic Characteristic Class]
Define
\[
\mathfrak{c}_{S}:=[dS]\in H^{1}(M\setminus\Sigma_{\mathrm{sing}};\mathbb{R}),
\]
viewed as a de~Rham cohomology class on the nonsingular locus.
\end{definition}

When $\Sigma_{\mathrm{sing}}$ has codimension two or higher, $\mathfrak c_{S}$
extends to a \emph{relative} cohomology class
\[
\mathfrak{c}_{S}\in H^{1}(M,\Sigma_{\mathrm{sing}};\mathbb{R}).
\]

\begin{proposition}[Flux Quantization (conditional)]
\label{prop:flux-quant}
If $\Sigma_{\mathrm{sing}}$ consists of isolated points and
the derived obstruction space has integer lattice structure
(e.g.\ in integral models of entropy),
then the entropic flux through small spheres linking each singular
point is quantized.
\end{proposition}

\begin{proof}[Proof (Sketch)]
Reduction to integral cohomology of the relative pair
$(M,\Sigma_{\mathrm{sing}})$; quantization then follows from integer
obstruction classes.
\end{proof}

\begin{remark}
This resembles magnetic charge quantization but arises purely from
entropy geometry.
\end{remark}



\section{Derived Resolution of Entropic Defects}
\label{sec:resolution}

Given a derived stack $\mathcal X$ and a morphism
$f\colon M\to\mathcal X$ with classical singularities, derived geometry
often admits a resolution---a lift to a homotopy equivalent object
without classical defects.

\begin{theorem}[Derived Resolution]
\label{thm:resolution}
Let $\Sigma_{\mathrm{sing}}$ be a closed set of entropic
singularities of $U$.
Then there exists a derived lift
\[
\widetilde U\colon M\longrightarrow\mathrm{Plenum}
\]
which is a derived resolution of $U$ in the sense that
$\widetilde U$ is a homotopy equivalence on the nonsingular locus and
extends uniquely across $\Sigma_{\mathrm{sing}}$ in the derived
category.
\end{theorem}

\begin{proof}[Proof (Sketch)]
This is a direct consequence of how derived mapping stacks encode
obstruction theory: the obstruction to extending $U$ lives in the
cotangent complex of $\mathrm{Plenum}$, which, by Theorem~\ref{thm:Main-Plenum},
has amplitude $[-1,1]$, forcing obstructions to vanish in strictly
positive degrees.
\end{proof}



\section{Physical Interpretation}
\label{sec:phys-top}

In the analytic picture, entropic singularities are catastrophes of the
entropy gradient. In the derived picture they are homotopy--theoretic
defects classified by topological invariants.  Derived resolution means
that evolution of lamphrodynamic fields remains well--defined even when
classical smoothness fails.  This is consonant with the BV formalism
(Ch.~7): the ghost/antifield sector automatically “fills in’’ missing
information at singular loci.



\section{Summary and Forward Link}
\label{sec:summary12}

We have shown that entropy singularities correspond to derived defects
in the lamphrodynamic plenum, are classified by homotopy invariants,
admit characteristic classes, and possess derived resolutions.
These ideas form the foundation for the geometric reductions of
Chapter~\ref{ch:symmetries-reduction}.


\chapter{Symmetries and Reduction}
\label{ch:symmetries-reduction}

Classical lamphrodynamics features a hierarchy of symmetries acting on
the fields $(\Phi,\vec v,S)$: internal symmetries, spacetime
diffeomorphisms, and entropy--preserving reparametrizations.
At the derived level these are encoded as group actions on the plenum
stack.  Reduction of shifted symplectic structures along such actions
produces reduced phase spaces and determines the structure of the BV
complex.  In this chapter we formulate these ideas rigorously, develop
quotient stacks and reduction, and establish the geometric foundations
for BV--BFV boundary theory.



\section{Group Actions on the Plenum Stack}
\label{sec:group-actions}

Let $G$ be a (possibly infinite--dimensional) Lie group acting smoothly
on the configuration bundle of lamphrodynamic fields.
Passing to the derived stack level,
\[
G\curvearrowright\mathrm{Plenum}.
\]

\begin{definition}[Action on Derived Stacks]
A \emph{derived action} of $G$ on $\mathrm{Plenum}$
is a morphism of stacks
\[
\alpha\colon B G\times \mathrm{Plenum} \longrightarrow \mathrm{Plenum}
\]
compatible with group multiplication in the homotopical sense.
\end{definition}

\begin{remark}
The formalism of $BG$ allows the action to carry higher homotopies, so
derivations of the action may vanish only up to coherent homotopy.
\end{remark}



\section{Quotient Stacks and Derived Reduction}
\label{sec:quotients}

We recall the standard homotopy quotient construction.

\begin{definition}[Homotopy Quotient]
The \emph{derived quotient stack} is
\[
\mathrm{Plenum}\sslash G :=
\mathrm{Plenum}\times_{G} EG,
\]
where $EG\to BG$ is a contractible principal $G$--bundle.
\end{definition}

\begin{proposition}[Properties]
\label{prop:quotient-props}
The quotient stack $\mathrm{Plenum}\sslash G$
is a derived Artin stack with tangent complex of amplitude $[-1,1]$.
\end{proposition}

\begin{proof}[Proof (Sketch)]
The homotopy quotient preserves Artin representability and tangent
amplitude under mild hypotheses, using standard results in derived
geometry \cite{ToenVezzosi-HAGII,Calaque16}.
\end{proof}



\section{$(-1)$--Shifted Symplectic Reduction}
\label{sec:shifted-reduction}

Let $(\mathrm{Plenum},\omega)$ carry the $2$--shifted symplectic
structure constructed in Chapter~\ref{ch:plenum-construction}.
The quotient by a symmetry group $G$ inherits a shifted symplectic form.

\begin{theorem}[Shifted Symplectic Reduction]
\label{thm:shifted-reduction}
If $G$ acts on $\mathrm{Plenum}$ preserving $\omega$, then the
homotopy quotient $\mathrm{Plenum}\sslash G$ carries a
$(2-1)$--shifted symplectic structure, i.e.\ a $1$--shifted structure.
\end{theorem}

\begin{proof}[Proof (Sketch)]
This is a special case of shifted symplectic reduction
\cite{Calaque16}.  Preservation of $\omega$ ensures that
the descent to the quotient kills one degree of the shift.
\end{proof}

\begin{remark}
Under AKSZ descent on $\mathsf{T}[1]M$, the $1$--shift then produces a
$(-2)$--shift, which contributes to the ghost/antifield balance of
the BV theory.  Thus symmetry reduction influences the field content
after quantization.
\end{remark}



\section{Reduced Phase Spaces}
\label{sec:reduced-phase}

The classical (degree--$0$) phase space of lamphrodynamics is
\[
\mathcal{P}
:= \mathrm{Map}(M,\mathrm{Plenum})^{(0)},
\]
and the reduced phase space is obtained by quotienting by the derived
symmetry.

\begin{definition}[Reduced Classical Phase Space]
\[
\mathcal{P}_{\mathrm{red}}
:= \mathrm{Map}(M,\mathrm{Plenum}\sslash G)^{(0)}.
\]
\end{definition}

\begin{proposition}[Reduced Symplectic Structure]
\label{prop:reduced-symplectic}
$\mathcal{P}_{\mathrm{red}}$ inherits a classical symplectic form which
is the reduction of the classical form on $\mathcal{P}$.
\end{proposition}

\begin{proof}[Proof (Sketch)]
Restrict the shifted reduction of Theorem~\ref{thm:shifted-reduction}
to degree $0$.  Standard arguments in AKSZ theory apply.
\end{proof}



\section{Gauge Fixing and Derived Geometry}
\label{sec:gauge-fix}

Gauge fixing can be regarded as selecting a section of the projection
$\mathrm{Plenum}\to\mathrm{Plenum}\sslash G$.  Derived geometry provides
a canonical method of choosing such sections via homotopy transfer.

\begin{theorem}[Homotopy Gauge Fixing]
\label{thm:gauge-fix}
Under mild hypotheses on $\mathrm{Plenum}$, there exists a homotopy
splitting of the projection
\[
\mathrm{Plenum}\longrightarrow\mathrm{Plenum}\sslash G,
\]
so any gauge choice corresponds to a homotopy section.
\end{theorem}

\begin{proof}[Proof (Sketch)]
The projection is a homotopy epimorphism in the derived category; a
homotopy section exists by standard model category arguments.
\end{proof}

\begin{remark}
Homotopy gauge fixing means no physical content depends on gauge
choices; ghost/antifield sectors encode the homotopy redundancies
automatically.
\end{remark}



\section{BV--BFV Boundary Structure}
\label{sec:bv-bfv}

Boundary terms of the AKSZ action produce a BV--BFV (Batalin--Vilkovisky /
Batalin--Fradkin--Vilkovisky) theory on the boundary~$\partial M$,
encoding the symplectic structure of boundary degrees of freedom.

\begin{theorem}[BV--BFV Boundary Theorem]
\label{thm:bv-bfv}
If $M$ has nonempty boundary, the AKSZ--RSVP action induces a
$(-1)$--shifted symplectic form on the boundary field space and a
BV--BFV structure consistent with the reduced phase space of the
bulk theory.
\end{theorem}

\begin{proof}[Proof (Sketch)]
This follows from the general AKSZ--BV--BFV correspondence
\cite{CattaneoFelder2012}, applied to the $2$--shifted plenum and its
boundary reduction.
\end{proof}



\section{Summary and Forward Link}
\label{sec:summary13}

We have developed derived symmetries, quotient stacks, shifted
symplectic reduction, and reduced phase spaces for lamphrodynamic
fields.  Gauge fixing is homotopy--theoretic, and BV--BFV boundary
structures appear automatically.  These geometric foundations prepare
for the detailed classical comparisons of the next chapter.



\chapter{Embedding Classical Field Theories into the Lamphrodynamic Plenum}
\label{ch:embedding}

\subsection{Introduction}

Classical field theory is usually formulated in terms of sections of bundles over a fixed spacetime manifold, together with variational principles determining the Euler--Lagrange equations. In the present volume we have recast field theory as arising from derived moduli spaces equipped with a $(-1)$--shifted symplectic structure. The goal of this chapter is to relate the ordinary formulation to the lamphrodynamic formalism by constructing a functor from a suitable category of classical field theories to the category representing RSVP field models.

We adopt the interpretative stance that classical theories embed \emph{functorially} into RSVP. That is, there exists an $(\infty,1)$--functor
[
F \colon \CFTh \longrightarrow \RSVP,
]
which is fully faithful (or at least conservative), and whose essential image contains precisely those lamphrodynamic models that arise from standard classical physics. In particular, the familiar theories of scalar fields, gauge fields, and Einstein gravity all appear as objects inside RSVP rather than as limits or degenerations.

\begin{remark}
This point of view corresponds to Option~C of the discussion ending the previous chapter: RSVP is a categorical enlargement of classical theory, not merely a generalization nor an approximation scheme.
\end{remark}

\subsection{The category of classical field theories}

Let $X$ be a smooth oriented spacetime manifold and let $E \to X$ be a smooth vector bundle (or principal bundle). A classical field configuration is a smooth section $\phi \colon X \to E$. A classical Lagrangian density $\mathcal{L}$ is a smooth function of the jet bundle $J^\infty(E)$, and the classical action is
[
S_{\mathrm{class}}[\phi] = \int_X \mathcal{L}(\phi, \nabla\phi, \ldots ), .
]

\begin{definition}
Define the $(\infty,1)$--category $\CFTh$ whose objects are triples $(X,E,\mathcal{L})$ and whose morphisms are field--bundle--variational morphisms preserving the Lagrangian structure.
\end{definition}

This category includes scalar field theory, Yang--Mills theory, and Einstein gravity, among others. It constitutes the standard domain of classical theoretical physics.

\subsection{Construction of the embedding functor}

We now define a functor
[
F \colon \CFTh \longrightarrow \RSVP.
]
The idea is that a classical field $\phi \colon X \to E$ is assigned a lamphrodynamic triple $(\Phi,\mathbf v,S)$ via a canonical lift. The potential $\Phi$ encodes the classical field value, while $\mathbf v$ and $S$ are obtained from trivial (or canonical) background data determined by the geometry of $X$ and the structure of $E$. Formally, set
[
F(X,E,\mathcal{L}) = \Map(X,\Plenum_{\Phi,\mathbf v,S}),
]
where $\Plenum_{\Phi,\mathbf v,S}$ is the lamphrodynamic moduli stack introduced previously, equipped with the $(-1)$--shifted symplectic structure constructed in Chapter~6.

\begin{proposition}
The assignment $F$ extends to an $(\infty,1)$--functor $F \colon \CFTh \to \RSVP$, preserving homotopy equivalences and sending classical Lagrangian morphisms to lamphrodynamic morphisms.
\end{proposition}

\begin{proof}[Sketch of proof]
A morphism of classical field theories induces a morphism of derived mapping stacks. The shifted symplectic forms are natural with respect to these maps by functoriality of the AKSZ construction. Hence the construction is functorial. The full proof uses the theory of tangent complexes discussed in Chapter~2.
\end{proof}

\subsection{Examples}

\subsubsection{Scalar fields}

Let $(X,E,\mathcal{L})$ be a scalar field theory with $E = X \times \mathbb{R}$ and
[
\mathcal{L} = \frac12(\partial\phi)^2 - V(\phi).
]
Then
[
F(X,E,\mathcal{L}) = \Map(X,\Plenum),
]
with $F(\phi) = (\Phi,\mathbf v,S)$ given by $\Phi(x) = \phi(x)$ and $\mathbf v,S$ trivial. The Euler--Lagrange equations become the $\Phi$--component of the lamphrodynamic equations.

\subsubsection{Yang--Mills fields}

For Yang--Mills theory the classical gauge potential determines a component of $\Phi$ via a fixed embedding $\mathfrak{g}\hookrightarrow\text{scalar sector}$, while $\mathbf v$ and $S$ remain trivial. The shifted symplectic form lifts from the moduli of gauge fields by AKSZ.

\subsubsection{Einstein gravity}

In the case of Einstein gravity, the Lorentzian metric determines $\Phi$ through a chosen geometric invariant such as curvature, and the components $\mathbf v$ and $S$ encode kinematic and entropic contributions from spacetime geometry. Details require the constructions of Chapter~12.

\subsection{Conclusion}

The functor $F$ embeds classical field theory into the RSVP formalism as a distinguished substructure. In this manner, the classical paradigm appears as a special case of a more general lamphrodynamic framework, and quantization can proceed categorically at the level of RSVP rather than theory by theory.

\begin{remark}
This approach allows one to regard classical field theories and RSVP models within a unified mathematical language, and it clarifies the relationship between classical physics and the derived geometric structures developed in the preceding chapters.
\end{remark}



% ------------------------------------------------------------
% Appendices
% ------------------------------------------------------------
\appendix

\appendix

\section*{Appendix A: Review of Differential Geometry}
\addcontentsline{toc}{section}{Appendix A: Review of Differential Geometry}

\section{Manifolds and Smooth Maps}

A \emph{smooth manifold} of dimension $n$ is a second countable Hausdorff space $M$ together with an atlas of coordinate charts ${(U_i,\varphi_i)}$ where $\varphi_i\colon U_i\to\mathbb{R}^n$ are homeomorphisms onto their images and all transition maps $\varphi_j\circ\varphi_i^{-1}$ are smooth.

\begin{definition}
A map $f\colon M\to N$ between smooth manifolds is called smooth if for every pair of charts $\varphi$ on $M$ and $\psi$ on $N$, the coordinate representation $\psi\circ f\circ\varphi^{-1}$ is a smooth function between Euclidean spaces.
\end{definition}

Note that in the categorical viewpoint, manifolds form a subcategory of smooth schemes, and ultimately of derived stacks.  The lamphrodynamic formalism requires a broad generalization of these notions, but ordinary smooth manifolds remain the geometric base of physical spacetime throughout Volume~I.

\subsection{Tangent Bundle}

For each $p\in M$, the tangent space $T_pM$ is defined as equivalence classes of smooth curves passing through $p$, or equivalently as derivations on germs of smooth functions at $p$.  The union of all tangent spaces
[
TM=\bigsqcup_{p\in M}T_pM
]
is a smooth vector bundle over $M$.

\begin{remark}
Later, tangent complexes of derived stacks generalize these tangent spaces.  In the derived setting, the tangent object is a chain complex rather than an ordinary vector space.
\end{remark}

\section{Vector Bundles and Differential Forms}

A (real) vector bundle over $M$ is a smooth manifold $E$ together with a smooth projection $\pi\colon E\to M$ such that for each $p\in M$, the fiber $E_p=\pi^{-1}(p)$ is a real vector space and locally $E\simeq U\times\mathbb{R}^k$.

\begin{definition}
The space of smooth sections of $E$ is denoted $\Gamma(M,E)$.
\end{definition}

In particular, the cotangent bundle $T^*M$ yields the spaces of differential $k$-forms
[
\Omega^k(M)=\Gamma(M,\wedge^kT^*M).
]

\section{Connections and Curvature}

Let $E\to M$ be a vector bundle. A connection is a $\mathbb{R}$--linear operator
[
\nabla\colon \Gamma(M,E)\to\Gamma(M,T^*M\otimes E)
]
satisfying the Leibniz rule $\nabla(fs)=df\otimes s + f,\nabla s$ for $f\in C^\infty(M)$ and $s\in\Gamma(M,E)$.

\begin{definition}
The curvature of a connection $\nabla$ is the bundle map
[
R_\nabla\colon \Gamma(M,E)\to \Gamma(M,\wedge^2T^*M\otimes E)
]
defined by $R_\nabla(s)=\nabla^2s$.
\end{definition}

Connections appear in RSVP as part of geometric background data, but the lamphrodynamic fields themselves are defined intrinsically from the derived moduli stack, not imposed externally.

\section{Riemannian and Lorentzian Geometry}

A Riemannian metric $g$ on $M$ is a smooth section of $T^*M\otimes T^*M$ which is symmetric and positive definite.  A Lorentzian metric has signature $(-,+,\ldots,+)$ and is the standard setting for general relativity.

The Levi--Civita connection is the unique torsion--free connection compatible with a given metric.  Its curvature tensor and associated invariants (Ricci curvature, scalar curvature) appear throughout classical field theory and reappear in the comparison sections of Chapter~14.

\begin{remark}
One conceptual theme of RSVP is that the lamphrodynamic metric need not be an externally fixed Lorentzian structure. Instead, the triple $(\Phi,\mathbf v,S)$ plays the role of the fundamental geometric data.
\end{remark}

\section{Exterior Calculus}

The exterior derivative $d:\Omega^k(M)\to\Omega^{k+1}(M)$ satisfies $d^2=0$, and the wedge product makes $\Omega^\bullet(M)$ into a graded commutative algebra.  Many constructions in the $(-1)$--shifted symplectic setting rely on derived analogues of differential forms.

\begin{definition}
The de,Rham cohomology of $M$ is the cohomology of the complex $(\Omega^\bullet(M),d)$.
\end{definition}

\begin{remark}
For derived stacks, the de,Rham complex is upgraded to a mixed complex whose cohomology yields the shifted symplectic form used in the AKSZ construction.
\end{remark}

\section{End of Appendix A}

This appendix only summarizes the necessary smooth differential geometry. The derived generalizations in Chapters~2--4 systematically replace tangent spaces by tangent complexes, replace manifolds by derived Artin stacks, and replace ordinary forms by their graded, shifted, or mixed counterparts.


\section*{Appendix B: Homological Algebra Essentials}
\addcontentsline{toc}{section}{Appendix B: Homological Algebra Essentials}

\section{Chain Complexes}

Let $\mathbf{Vect}$ denote the category of real vector spaces. A chain complex $(C_\bullet,d)$ is a sequence of vector spaces
[
\cdots \longrightarrow C_{n+1}\xrightarrow{d_{n+1}} C_n\xrightarrow{d_n} C_{n-1}\longrightarrow\cdots
]
such that $d_n\circ d_{n+1}=0$ for all $n\in\mathbb{Z}$.  One defines the $n^{\mathrm{th}}$ homology as
[
H_n(C_\bullet)=\frac{\ker(d_n)}{\operatorname{im}(d_{n+1})}.
]

\begin{remark}
Tangent complexes of derived stacks are examples of chain complexes whose homology measures infinitesimal automorphisms and obstructions.
\end{remark}

\section{Double Complexes and Totalization}

A double complex $C_{\bullet,\bullet}$ is a bi-graded collection ${C_{p,q}}$ with horizontal and vertical differentials. The total complex $\operatorname{Tot}(C)$ is defined by
[
\operatorname{Tot}(C)*n = \bigoplus*{p+q=n} C_{p,q}.
]
Totalization is crucial in defining the cotangent complex and for bookkeeping in BV--BRST constructions.

\section{Tensor Products and Hom Functors}

For chain complexes $C_\bullet$ and $D_\bullet$, the tensor product is the complex
[
(C\otimes D)*n = \bigoplus*{p+q=n} C_p\otimes D_q,
]
with differential determined by the Leibniz rule.  Dually, the internal Hom complex $\underline{\mathrm{Hom}}(C,D)$ is given by
[
\underline{\mathrm{Hom}}(C,D)*n = \prod_m \mathrm{Hom}(C_m,D*{m+n}),
]
with the standard differential. This is the backbone of enriched homotopical structure.

\section{Exact Sequences and Derived Functors}

A short exact sequence of complexes
[
0\to A_\bullet\to B_\bullet\to C_\bullet\to0
]
induces a long exact homology sequence. Derived functors such as $\operatorname{Ext}$ and $\operatorname{Tor}$ arise by taking projective or injective resolutions of complexes.

\begin{remark}
For a derived stack $\mathcal X$, its cotangent complex $\mathbb{L}*\mathcal X$ and tangent complex $\mathbb{T}*\mathcal X$ are defined using homological algebra in suitable model categories.
\end{remark}

\section{Model Categories}

A model category is a category equipped with three distinguished classes of morphisms (weak equivalences, fibrations, cofibrations) satisfying appropriate axioms. Homotopy categories and derived functors are constructed by formally inverting weak equivalences.

\begin{definition}
A Quillen adjunction between model categories induces an adjunction on homotopy categories, and a Quillen equivalence induces an equivalence of homotopy categories.
\end{definition}

Model category language underlies much of the theory of derived algebraic geometry and is assumed in the main text.  The reader seeking more detail will find standard references in the bibliography.

\section{Spectral Sequences}

Spectral sequences are computational tools associated to filtered chain complexes. They provide successive approximations to homology and cohomology and are indispensable for computing deformation and obstruction theories in derived geometry.

\begin{remark}
The obstruction theory used to construct the moduli stack $\Plenum$ depends on the vanishing of certain classes in higher derived $\operatorname{Ext}$ spaces, often computed via spectral sequence arguments.
\end{remark}

\section{End of Appendix B}

We emphasize that the natural home for RSVP is the $(\infty,1)$--categorical world introduced in Chapter~2, in which many of the above constructions are given a homotopical rather than merely algebraic interpretation.


\section*{Appendix C: Functional Analysis for PDEs}
\addcontentsline{toc}{section}{Appendix C: Functional Analysis for PDEs}

\section{Function Spaces}

Let $\Omega\subset\mathbb{R}^n$ be an open set. The space $L^2(\Omega)$ consists of square--integrable functions with norm
[
|u|*{L^2}^2 = \int*\Omega |u(x)|^2,dx.
]
More generally, $L^p(\Omega)$ denotes the space of $p$--integrable functions.  The Sobolev space $H^k(\Omega)$ is defined by
[
H^k(\Omega)={ u\in L^2(\Omega)\mid \partial^\alpha u\in L^2(\Omega)\text{ for all }|\alpha|\le k}.
]

\begin{remark}
Sobolev spaces provide the natural framework for weak solutions of lamphrodynamic equations and for proving well--posedness theorems.
\end{remark}

\section{Weak Derivatives and Weak Solutions}

Let $u\in L^1_{\mathrm{loc}}(\Omega)$. A weak derivative $\partial_i u$ is a distribution satisfying
[
\int_\Omega u,\partial_i\varphi=-\int_\Omega (\partial_i u),\varphi
]
for all test functions $\varphi\in C_c^\infty(\Omega)$. A weak solution of a PDE is a function satisfying the PDE in the sense of distributions.

\begin{definition}
A function $u\in H^1(\Omega)$ is called a weak solution to a differential operator $L$ if
[
\langle Lu,\varphi\rangle=0
]
for all test functions $\varphi$.
\end{definition}

\section{Energy Methods}

Consider an evolution equation
[
\partial_t u +Lu = N(u),
]
where $L$ is a linear operator and $N$ a nonlinear term. Multiplying by $u$ and integrating by parts produces an energy identity of the form
[
\frac{d}{dt}|u|_{L^2}^2 + \langle Lu, u\rangle = \langle N(u),u\rangle.
]
Such estimates are central to the lamphrodynamic theory discussed in Chapters~9--11.

\begin{remark}
The entropy field $S$ introduces additional dissipation, strengthening the coercivity of energy estimates under suitable assumptions.
\end{remark}

\section{Semigroups and Linear Evolution}

For linear equations $\partial_t u+Lu=0$ on a Banach space $X$, well--posedness is often obtained by showing that $-L$ generates a strongly continuous semigroup ${e^{-tL}}_{t\ge0}$ on $X$.

\begin{theorem}[Hille--Yosida]
Let $X$ be a Banach space.  A densely defined closed operator $-L$ generates a contraction semigroup on $X$ if and only if its resolvent satisfies suitable positivity and growth conditions.
\end{theorem}

This theorem is used in the proof of global well--posedness for linearized lamphrodynamic fields.

\section{Compactness and Regularity}

Weak compactness and Rellich--Kondrachov compactness theorems allow one to extract convergent subsequences in Sobolev spaces. Elliptic regularity for elliptic operators $L$ improves the regularity of weak solutions:
[
Lu=f\in H^k\quad\Rightarrow\quad u\in H^{k+2}.
]

\section{Nonlinear PDE Techniques}

Local existence for nonlinear systems often follows from fixed point arguments (Picard iteration, contraction mapping principle) in suitable function spaces. Global existence usually requires energy estimates and a priori bounds.

\begin{remark}
Blow--up phenomena are typically associated with failure of energy control or growth of nonlinear terms. Entropy production in RSVP plays a role in mitigating such instabilities.
\end{remark}

\section{End of Appendix C}

This appendix summarizes the analytic tools required for the proofs in Part~III of this volume.  The reader is referred to standard texts on PDEs and functional analysis for details beyond this short review.


\section*{Appendix D: Conventions and Notation}
\addcontentsline{toc}{section}{Appendix D: Conventions and Notation}

\section{Indices and Summation}

Throughout this monograph, lower case Latin indices $i,j,k,\dots$ range over spatial components $1,\dots,n$, while Greek indices $\mu,\nu,\rho,\sigma$ range over spacetime indices $0,\dots,n$.  The Einstein summation convention over repeated indices is always assumed unless stated otherwise.

Boldface symbols such as $\mathbf v$ denote spatial vector fields, whereas non-bold symbols $v^\mu$ denote spacetime components.

\section{Metrics and Signatures}

Spacetime metrics $g_{\mu\nu}$ have Lorentzian signature $(-,+,\ldots,+)$ unless otherwise stated.  Raising and lowering of indices is performed with $g_{\mu\nu}$ and its inverse $g^{\mu\nu}$.

The metric used in classical background comparisons may differ from the effective geometric structures arising in lamphrodynamics; the latter are derived from the triple $(\Phi,\mathbf v,S)$ rather than postulated \emph{a priori}.

\section{Differential Operators}

The gradient operator $\nabla$ denotes the Levi--Civita covariant derivative unless the context dictates a gauge or derived connection. The d'Alembertian operator is
[
\square = \nabla^\mu\nabla_\mu.
]

The divergence of a vector field is $\nabla\cdot\mathbf v$, and the Laplacian on scalar functions is $\Delta=\nabla^i\nabla_i$.  In derived settings, the operator $\Delta$ refers to the BV Laplacian unless explicitly disambiguated.

\section{Lamphrodynamic Fields}

The symbol $\Phi$ denotes the scalar potential field.  The bold symbol $\mathbf v$ denotes the spatial vector field of lamphrodynamic flow, and $S$ denotes the entropy density.  The triple $(\Phi,\mathbf v,S)$ is the fundamental object of RSVP field theory.

\begin{remark}
It is important not to confuse the entropy field $S$ with the classical thermodynamic entropy, although there are interpretative analogies.  In RSVP, $S$ is a geometric quantity arising from derived symplectic structures.
\end{remark}

\section{Derived Geometry}

The symbols $\mathbb{L}*\mathcal X$ and $\mathbb{T}*\mathcal X$ denote the cotangent and tangent complexes of a derived stack $\mathcal X$.  The shifted symplectic structure on $\Plenum$ is denoted by $\omega_{[-1]}$.

\section{Variational Notation}

The classical BV action is denoted $\mathcal{S}$, and the classical master equation is ${\mathcal S,\mathcal S}=0$.  The BV bracket is written ${-,-}_{BV}$ or simply ${-,-}$ when unambiguous.

\section{Constants and Parameters}

Physical constants such as $c$ and $G$ may be set to unity except in comparative formulas with observational cosmology.  Dimensional analysis is included in Volume~II as needed.

\section{End of Appendix D}

These conventions are used uniformly throughout all volumes of the monograph.  Any deviation in later chapters will be indicated explicitly.


\section*{Appendix E: Proof of Technical Lemmas}
\addcontentsline{toc}{section}{Appendix E: Proof of Technical Lemmas}

\section{Introduction}

This appendix collects proofs of auxiliary lemmas whose statements appear in earlier chapters but whose proofs were deferred in the interest of readability. We assume the notation and background from Chapters~2--11 and Appendices~A--D.

Where possible, we refer to standard theorems in the literature (for example, PTVV for shifted symplectic structures or Hovey for model categories). Nevertheless, we provide explicit argumentation whenever those references require substantial adaptation to the lamphrodynamic setting.

\section{Lemma 3.4: Exactness of the Shifted de,Rham Differential}

\begin{lemma}[Lemma 3.4]
Let $\omega_{[-1]}$ be the $(-1)$--shifted presymplectic form on a derived Artin stack $\mathcal{X}$. Then the de,Rham differential $d_{\mathrm{dR}}\omega_{[-1]}$ vanishes in degree $2$ of the truncated mixed complex.
\end{lemma}

\begin{proof}
By PTVV, a $k$--shifted symplectic structure is a closed 2--form in the truncated de,Rham complex of degree $k$. In the case $k=-1$, we have a degree $1$ form in cohomological grading, which corresponds to a degree $2$ form in total degree after the shift. The closure condition then asserts $d_{\mathrm{dR}}\omega_{[-1]}=0$. Since the truncation functor commutes with $d_{\mathrm{dR}}$ on derived Artin stacks, the vanishing holds in the truncated complex as required.
\end{proof}

\section{Lemma 6.2: Smoothness Condition for the Moduli Stack $\Plenum$}

\begin{lemma}[Lemma 6.2]
Let $\Plenum$ denote the lamphrodynamic moduli stack defined in Chapter~6. Then $\Plenum$ is a derived Artin stack locally of finite presentation.
\end{lemma}

\begin{proof}[Sketch of proof]
Recall that $\Plenum$ is constructed as a derived mapping stack $\Map(X,\mathcal{F})$, where $\mathcal{F}$ is a derived stack parameterizing the lamphrodynamic triple $(\Phi,\mathbf v,S)$. Since $X$ is smooth and $\mathcal{F}$ is locally of finite presentation, standard results on derived mapping stacks imply that $\Map(X,\mathcal{F})$ is a derived Artin stack locally of finite presentation. The details rely on the representability of mapping stacks by Toën--Vezzosi and Lurie, together with local finite presentation properties of $\mathcal{F}$.
\end{proof}

\section{Lemma 9.7: Sobolev Inequality in Lamphrodynamic Variables}

\begin{lemma}[Lemma 9.7]
Let $\Phi,\mathbf v,S$ satisfy the hypotheses of Chapter~9. Then there exists a constant $C>0$ such that
[
|\Phi|*{L^\infty(\Omega)} \le C|\Phi|*{H^k(\Omega)}
]
for any open bounded domain $\Omega\subset\mathbb{R}^n$ and integer $k>\frac{n}{2}$.
\end{lemma}

\begin{proof}
This is the classical Sobolev embedding theorem. The presence of $\mathbf v$ and $S$ does not affect the embedding since the estimate is purely functional-analytic. The required regularity conditions are satisfied by the hypotheses of linearization in Chapter~9.
\end{proof}

\section{Lemma 10.5: A Priori Energy Estimate}

\begin{lemma}[Lemma 10.5]
Under the assumptions of Theorem~10.3, the a priori energy estimate
[
|(\Phi,\mathbf v,S)(t)|*{H^1}^2 \le C\exp\left(CT\right)|(\Phi,\mathbf v,S)(0)|*{H^1}^2
]
holds for all $t\in[0,T]$.
\end{lemma}

\begin{proof}
Multiply the linearized lamphrodynamic equations by $(\Phi,\mathbf v,S)$ and integrate in space. Integration by parts eliminates first--order derivatives. The entropy term contributes a negative definite form, and Grönwall's inequality yields the stated estimate.
\end{proof}

\section{Lemma 11.8: Local Existence for the Nonlinear System}

\begin{lemma}[Lemma 11.8]
Assume the nonlinear lamphrodynamic system satisfies the hypotheses of Chapter~11. Then there exists a time $T>0$ such that a unique solution exists on $[0,T]$ in $H^k$ for sufficiently large $k$.
\end{lemma}

\begin{proof}
Local existence follows from a fixed--point argument (Picard iteration) in $H^k$ spaces. The nonlinear terms satisfy a Lipschitz condition in $H^k$ due to Sobolev embedding and Moser estimates. Uniqueness follows from a Grönwall argument applied to the difference of two solutions.
\end{proof}

\section{End of Appendix E}

\begin{center}
\textit{This concludes Volume~I: Mathematical Foundations of the Relativistic Scalar--Vector Plenum.}
\end{center}



\backmatter
\printbibliography

\end{document}

